\section{What to Automate: A Task-Level Model}\label{sec:tasks}\label{sec:core}
% Generated by code/examples/invoice_chain.py from data/examples/invoice_chain.yaml; do not edit.
% Every number of the running example in part2/theory_core.tex is one of these macros.
\newcommand{\icVolume}{120{,}000}
\newcommand{\icRevenue}{\$1.2 billion}
\newcommand{\icWage}{\$34.80}
\newcommand{\icWageAnnual}{\$50{,}670}
\newcommand{\icIotaLow}{0.4}
\newcommand{\icIotaHigh}{0.9}
\newcommand{\icIotaN}{0.8}
\newcommand{\icR}{8.93\%}
\newcommand{\icTA}{4}
\newcommand{\icTh}{4}
\newcommand{\icThEngine}{10}
\newcommand{\icDR}{0.65}
\newcommand{\icDRlow}{0.45}
\newcommand{\icDRhigh}{0.85}
\newcommand{\icDh}{0.4}
\newcommand{\icRhoO}{65\%}
\newcommand{\icRhoRev}{0.575}
\newcommand{\icWagePremium}{1.2}
\newcommand{\icBetaU}{0.25}
\newcommand{\icEzero}{0.25}
\newcommand{\icEone}{0.033}
\newcommand{\icHzero}{12}
\newcommand{\icEUdecision}{0.5}
\newcommand{\icRowCapture}{Capture & 0.05 & 1.74 & 0.60, $N$ & 0.07 & 0.468 / 0.070 & 0.98 & 0.713 & 1.61 & 138}
\newcommand{\icLCapture}{\$17.40}
\newcommand{\icRowMatch}{Match & 0.08 & 2.78 & 2.17, $O$ & 0.11 & 0.183 / 0.077 & 0.98 & 0.741 & 2.52 & 161}
\newcommand{\icLMatch}{\$27.84}
\newcommand{\icRowExceptions}{Exceptions & 0.12 & 4.18 & 3.51, $U$ & 0.16 & 0.191 / 0.086 & 0.97 & 0.652 & 3.02 & 185}
\newcommand{\icLExceptions}{\$12.53}
\newcommand{\icRowApproval}{Approval & 0.05 & 1.74 & 1.55, $U$ & 0.07 & 0.177 / 0.070 & 0.99 & 0.857 & 1.61 & 138}
\newcommand{\icLApproval}{\$17.40}
\newcommand{\icLmult}{10, 10, 3 and 10}
\newcommand{\icEhighMin}{7.0\%}
\newcommand{\icEhighMax}{8.6\%}
\newcommand{\icDstarMin}{0.652}
\newcommand{\icDstarMax}{0.86}
\newcommand{\icMaxRelFullMin}{98.0\%}
\newcommand{\icMaxRelFullMax}{98.3\%}
\newcommand{\icMaxRelExcFull}{97.7\%}
\newcommand{\icEfullMin}{4.8\%}
\newcommand{\icEcapLow}{36\%}
\newcommand{\icApiReadiness}{0.6}
\newcommand{\icPinTotal}{\$0.40}
\newcommand{\icKappaO}{0.10}
\newcommand{\icZetaO}{0.03}
\newcommand{\icQMstep}{0.995}
\newcommand{\icOwnedOpex}{30\%}
\newcommand{\icDeltaKcompute}{25\%}
\newcommand{\icCapPerDollar}{3}
\newcommand{\icDipShare}{10\%}
\newcommand{\icCint}{\$3{,}000}
\newcommand{\icScaleCurv}{\$1{,}000}
\newcommand{\icQbarHigh}{98\%}
\newcommand{\icQbarExc}{97\%}
\newcommand{\icContractYears}{3}
\newcommand{\icExcIotaFeas}{0.902}
\newcommand{\icExcIotaStarDefault}{0.97}
\newcommand{\icDfeas}{0.86}
\newcommand{\icDill}{0.85}
\newcommand{\icQcap}{98\%}
\newcommand{\icQstar}{98.8\%}
\newcommand{\icResidCap}{2\%}
\newcommand{\icResidPolicy}{2\%}
\newcommand{\icQbarAppr}{99\%}
\newcommand{\icResidAppr}{1\%}
\newcommand{\icLmultReg}{10}
\newcommand{\icLmultRegLow}{5}
\newcommand{\icLmultRegHigh}{30}
\newcommand{\icRbarHighLow}{97.5\%}
\newcommand{\icRbarHighHigh}{99.5\%}
\newcommand{\icCaptureIotaFeas}{0.959}
\newcommand{\icMatchIotaFeas}{0.994}
\newcommand{\icMaxRelApprFull}{98.3\%}
\newcommand{\icAltWfull}{$-$\$14{,}000}
\newcommand{\icAltPair}{$-$\$1{,}500}
\newcommand{\icAltPairLoss}{\$1{,}500}
\newcommand{\icAltFirstIota}{0.903}
\newcommand{\icAltIotaStarN}{0.912}
\newcommand{\icRegMult}{1.5}
\newcommand{\icEHuman}{2\%}
\newcommand{\icQHuman}{98\%}
\newcommand{\icIotaStarP}{0.90}
\newcommand{\icExcAninety}{\$2.86}
\newcommand{\icXiMin}{0.77}
\newcommand{\icXiMax}{0.87}
\newcommand{\icMatchH}{\$2.78}
\newcommand{\icMatchO}{\$2.17}
\newcommand{\icMatchA}{\$2.52}
\newcommand{\icMatchVsH}{\$32{,}000}
\newcommand{\icMatchVsO}{\$42{,}000}
\newcommand{\icWinCzeroA}{\$2.41}
\newcommand{\icExcU}{\$3.51}
\newcommand{\icExcAone}{\$2.57}
\newcommand{\icExcAill}{\$3.02}
\newcommand{\icExcAeightyfive}{\$3.02}
\newcommand{\icExcDtie}{0.73}
\newcommand{\icExcDinf}{0.652}
\newcommand{\icCmanHigh}{\$1.05 million}
\newcommand{\icEinhouse}{\$1.25 million}
\newcommand{\icDivUnit}{$-$\$112{,}000}
\newcommand{\icDivCoord}{\$98{,}000}
\newcommand{\icDivRel}{$-$\$84{,}000}
\newcommand{\icDivOver}{\$93{,}000}
\newcommand{\icQM}{0.980}
\newcommand{\icQA}{0.913}
\newcommand{\icWfullHigh}{$-$\$5{,}500}
\newcommand{\icSingleCapture}{$-$\$168{,}000}
\newcommand{\icSingleMatch}{$-$\$116{,}000}
\newcommand{\icSingleExc}{$-$\$28{,}000}
\newcommand{\icSingleAppr}{$-$\$54{,}000}
\newcommand{\icBestPair}{\$6{,}600}
\newcommand{\icPairEA}{\$6{,}600}
\newcommand{\icPairEAloss}{$-$\$6{,}600}
\newcommand{\icBestFirstIota}{0.889}
\newcommand{\icIotaChainFeas}{0.65}
\newcommand{\icWeightyfive}{$-$\$60{,}000}
\newcommand{\icWninetyfive}{\$59{,}000}
\newcommand{\icWbelowN}{$-$\$18{,}000}
\newcommand{\icWaboveN}{$-$\$109{,}000}
\newcommand{\icWdropN}{\$91{,}000}
\newcommand{\icIotaStarN}{0.905}
\newcommand{\icXiPshare}{7.4\%}
\newcommand{\icLp}{\$10.44}
\newcommand{\icEtaMM}{\$0.31}
\newcommand{\icEtaAM}{\$0.54}
\newcommand{\icEtaAA}{\$0.39}
\newcommand{\icDetWdone}{\$175{,}000}
\newcommand{\icDetWdoneRun}{$-$\$154{,}000}
\newcommand{\icDetDWzero}{0.861}
\newcommand{\icW}{$-$\$164{,}000}
\newcommand{\icWbeforeRun}{\$165{,}000}
\newcommand{\icRunTotal}{\$329{,}000}
\newcommand{\icF}{\$2.81 million}
\newcommand{\icFbuild}{\$622{,}000}
\newcommand{\icFint}{\$2.00 million}
\newcommand{\icFother}{\$186{,}000}
\newcommand{\icIntCost}{\$60{,}000}
\newcommand{\icIntLow}{\$20{,}000}
\newcommand{\icIntHigh}{\$200{,}000}
\newcommand{\icNsys}{4}
\newcommand{\icCzero}{\$1.14 million}
\newcommand{\icCJ}{\$114{,}000}
\newcommand{\icCpar}{\$1.16 million}
\newcommand{\icCdup}{\$120{,}000}
\newcommand{\icCstr}{\$234{,}000}
\newcommand{\icCsurv}{\$0.00}
\newcommand{\icOs}{\$6{,}100}
\newcommand{\icNsv}{7.1}
\newcommand{\icExcessQuits}{2.35\%}
\newcommand{\icCrep}{0.56}
\newcommand{\icCwo}{\$300{,}000}
\newcommand{\icVarpi}{0.82}
\newcommand{\icBbar}{0.67}
\newcommand{\icTdet}{0.25}
\newcommand{\icLroll}{\$1.14 million}
\newcommand{\icOmegaPar}{80\%}
\newcommand{\icNuRec}{10\%}
\newcommand{\icTpar}{0.25}
\newcommand{\icTparStar}{0.00}
\newcommand{\icResid}{7\%}
\newcommand{\icNPV}{$-$\$3.78 million}
\newcommand{\icNPVcertain}{$-$\$3.97 million}
\newcommand{\icNPVunlimited}{$-$\$4.49 million}
\newcommand{\icNPVvarpiOne}{$-$\$3.79 million}
\newcommand{\icNPVthThree}{$-$\$3.71 million}
\newcommand{\icNPVthFive}{$-$\$3.84 million}
\newcommand{\icNPVthEngine}{$-$\$4.07 million}
\newcommand{\icNPVprogram}{$-$\$4.17 million}
\newcommand{\icRebuilds}{4 and 8}
\newcommand{\icFbreakeven}{$-$\$970{,}000}
\newcommand{\icNPVaLow}{$-$\$2.44 million}
\newcommand{\icNPVaDef}{$-$\$3.78 million}
\newcommand{\icNPVaHigh}{$-$\$8.44 million}
\newcommand{\icFintLow}{\$667{,}000}
\newcommand{\icFintHigh}{\$6.67 million}
\newcommand{\icNPVb}{$-$\$1.71 million}
\newcommand{\icNPVbCertain}{$-$\$1.90 million}
\newcommand{\icFb}{\$808{,}000}
\newcommand{\icSPcashOne}{0.32}
\newcommand{\icSPcashLater}{0.01}
\newcommand{\icSPgaapOne}{0.17}
\newcommand{\icSPgaapLater}{0.07}
\newcommand{\icYs}{$-$\$3.79 million}
\newcommand{\icYf}{$-$\$3.70 million}
\newcommand{\icYsb}{$-$\$1.72 million}
\newcommand{\icFpil}{\$143{,}000}
\newcommand{\icPilotShare}{15\%}
\newcommand{\icPilotFixed}{\$50{,}000}
\newcommand{\icDeltaM}{17\%}
\newcommand{\icTup}{1.22}
\newcommand{\icRevalShare}{10\%}
\newcommand{\icVolFactor}{2.01}
\newcommand{\icMaint}{\$106{,}000}
\newcommand{\icReval}{\$48{,}000}
\newcommand{\icTechDebt}{\$144{,}000}
\newcommand{\icNint}{16}
\newcommand{\icCcpl}{\$32{,}000}
\newcommand{\icCinc}{\$35{,}000}
\newcommand{\icSecSpend}{\$11{,}000}
\newcommand{\icSecResid}{\$20{,}000}
\newcommand{\icWrunHigh}{$-$\$335{,}000}
\newcommand{\icScope}{\$84{,}000}
\newcommand{\icCs}{\$20{,}000}
\newcommand{\icBBfix}{\$483{,}000}
\newcommand{\icBBbuilds}{\$185{,}000}
\newcommand{\icBBrun}{\$298{,}000}
\newcommand{\icBBpooled}{\$1.04}
\newcommand{\icBBint}{\$186{,}000}
\newcommand{\icBBfixv}{\$55{,}000}
\newcommand{\icBBcv}{\$1.41}
\newcommand{\icBBv}{304{,}000}
\newcommand{\icBBvnorun}{92{,}000}
\newcommand{\icBBsave}{\$259{,}000}
\newcommand{\icBBlock}{\$124{,}000}
\newcommand{\icBBsw}{\$373{,}000}
\newcommand{\icBBmarkup}{25\%}
\newcommand{\icScalePer}{\$154{,}000}
\newcommand{\icScaleFixed}{\$615{,}000}
\newcommand{\icScaleJensen}{\$20{,}000}
\newcommand{\icScalePool}{\$92{,}000}
\newcommand{\icCapDc}{\$0.28}
\newcommand{\icCapDK}{\$1.21}
\newcommand{\icCapG}{12\%}
\newcommand{\icCapValue}{\$0.13}
\newcommand{\icCapCash}{\$0.17}
\newcommand{\icCapGaap}{\$0.02}
\newcommand{\icCapPin}{\$0.40}
\newcommand{\icHalEo}{7.0\%}
\newcommand{\icHalEh}{0.28\%}
\newcommand{\icHret}{5\%}
\newcommand{\icHalFloor}{2.6\%}
\newcommand{\icHalFloorO}{2.4\%}
\newcommand{\icHalFloorDh}{2.6\%}
\newcommand{\icHalTol}{1\%}
\newcommand{\icEpsd}{0.5}
\newcommand{\icBiasPeak}{61\%}
\newcommand{\icBiasThr}{67\%}
\newcommand{\icBiasTop}{72\%}
\newcommand{\icBiasCentral}{55\%}
\newcommand{\icBiasUplift}{10\%}
\newcommand{\icAutL}{\$17.40}
\newcommand{\icAutLm}{\$30.45}
\newcommand{\icAutCap}{\$609}
\newcommand{\icAlpha}{1.5}
\newcommand{\icAlphaLow}{1.15}
\newcommand{\icAutCmTwo}{\$0.001}
\newcommand{\icAutCmThree}{\$0.004}
\newcommand{\icAutHumOne}{\$1.39}
\newcommand{\icAutHumTwo}{\$0.52}
\newcommand{\icAutCvar}{\$673}
\newcommand{\icAutCvarLow}{\$3{,}200}
\newcommand{\icAutCvarQ}{95\%}
\newcommand{\icAutLmax}{\$5.00 million}
\newcommand{\icCerrFloor}{\$0.45}
\newcommand{\icCerrMis}{\$0.004}
\newcommand{\icCerrAnnual}{\$55{,}000}
\newcommand{\icCerrNoRev}{\$1.27}
\newcommand{\icCerrNoRevAnnual}{\$152{,}000}
% running-example numbers, generated by code/examples/invoice_chain.py from data/examples/invoice_chain.yaml

Section~\ref{sec:model} takes the firm's cost structure $(\gamma,a,a_{\min},c_c,k)$ and its growth technology $\lambda$ as given. This section derives them from below. A firm is a set of activities organized into processes. Each activity can be done in house, outsourced, augmented by AI, automated or eliminated; the firm chooses modes, data-readiness investment and the timing of transitions; and the resulting cost structure is scored by Part I's maximal score $R_X^*$ and firm value. The unit of analysis is the task, as in task-based models of automation \citep{acemoglu2018race,restrepo2024automation}, but the comparison is between modes inside one firm's processes, and it includes the make-or-buy margin of transaction-cost and property-rights theories of the firm \citep{williamson1985economic,grossman1986costs}, which AI moves in both directions \citep{shahidi2025coasean,garicano2026vanishing}. Five conclusions stand out. The value of automation must be measured against the cheapest alternative, not against in-house labor, and outsourcing vendors that adopt AI can close the window for in-house automation (Section~\ref{sec:tasks-frontier}). Automation beats augmentation only when most units can skip human review, and only if reviewers catch enough of the AI's errors to meet a reliability target at all (Section~\ref{sec:tasks-review}). Handoffs and chained errors make whole processes the natural unit of automation, so that activity-by-activity business cases misstate its value (Section~\ref{sec:tasks-process}), and data readiness acts as a threshold rather than a dial (Section~\ref{sec:tasks-readiness}). Capital, run costs, security, AI failure modes, transition costs and execution risk change the ranking of projects, the scope of automation and the path of the Rule-of-40 score (Sections~\ref{sec:tasks-capital} to \ref{sec:tasks-failure}). Finally, every activity-level choice maps into the parameters of Section~\ref{sec:model}, and the automated share is a step function of capability and cost (Sections~\ref{sec:tasks-aggregation} and \ref{sec:tasks-breakpoints}). Proofs are in Appendix~\ref{app:proofs-tasks}. The script \texttt{code/verify\_part2\_core.py} checks every closed form, first-order condition and threshold, and every number in the examples.

\subsection{Activities, modes and unit costs}\label{sec:tasks-setup}

\begin{definition}[Activities and processes]\label{def:tasks-activity}
An activity $j$ is a detailed work activity performed in a given occupation (O*NET lists 2,087 detailed work activities; \citealp{onet2026database}). It has annual volume $v_j>0$ (units of output), in-house time $t_j>0$ (hours per unit) and loaded wage $w_j$ (dollars per hour, including payroll taxes and benefits). Its labor-equivalent spend is $E_j:=v_jw_jt_j$. A process $p$ is a set $\mathcal J_p$ of $J$ activities linked by directed handoff edges $\mathcal E_p$; edge $e$ carries $v_e$ units a year, the process delivers $v_p$ units of output, and its critical path $\mathcal C_p\subseteq\mathcal J_p$ is the set of activities every unit passes through.
\end{definition}

\begin{assumption}[Modes and unit costs]\label{ass:tasks-modes}
Activity $j$ can be done in house ($H$), outsourced ($O$), augmented ($U$), automated ($A$) or eliminated ($N$). Its feasible set $\mathcal M_j\ni H$ reflects capability, data readiness and permissibility. At an inference price $\pi$ (dollars per unit of inference: tokens or agent steps), the unit cost of mode $x$ is affine in $\pi$,
\begin{equation}\label{eq:tasks-cost}
c_j(x;\pi)=c^0_j(x)+\pi\,n_j(x),\qquad n_j(H)=n_j(O)=n_j(N)=0<n^U_j:=n_j(U)<n_j:=n_j(A),
\end{equation}
with non-inference costs $c^0_j(H)=w_jt_j$, $c^0_j(O)=p_{Oj}+\kappa_Ow_jt_j+\zeta_{Oj}$, $c^0_j(U)=w_jt^U_j+e^U_jL_j$, $c^0_j(A)=c^Q_j+F^A_j/(v_jT_A)$ and $c^0_j(N)=L^N_j$, each plus a capital charge $(r+\delta_K)K_j(x)$ (Section~\ref{sec:tasks-capital}). Here $p_{Oj}$ is the vendor's price, $\kappa_O$ the share of in-house cost spent on coordinating the vendor, $\zeta_{Oj}$ a risk and quality premium, $t^U_j\le t_j$ the augmented time per unit, $e^U_j$ the error rate that augmented work adds to in-house work (positive for tasks outside the AI's frontier, where people follow wrong suggestions, \citealp{dellacqua2026navigating}, and negative where assistance improves quality), $L_j$ the loss per error (Section~\ref{sec:tasks-review}), $c^Q_j$ the quality cost of automation (review plus undetected errors, Section~\ref{sec:tasks-review}), $F^A_j$ the build and integration cost amortized over $T_A$ years (with discounting, $1/T_A$ becomes the capital recovery factor $r/[1-\exp(-rT_A)]$, and every result below holds with either), $L^N_j$ the value lost per unit when a redesign eliminates the activity, and $K_j(x)$ the capital per unit of annual volume. Write $\underline c_j(\pi):=\min_{x\in\mathcal M_j\setminus\{A\}}c_j(x;\pi)$ for the best alternative to automation and $c^{\mathrm{pre}}_j:=\min_{x\in\mathcal M_j\cap\{H,O,N\}}c_j(x)$ for the best cost without AI. Each activity is flagged as a decision task (classifying, matching, approving) or a creation task (drafting, designing, coding), because the two respond differently to augmentation (Proposition~\ref{prop:tasks-augment}(d)).
\end{assumption}

Table~\ref{tab:tasks-notation} collects the notation. Part I's symbols keep their meanings.

\begin{table}[t]
\centering
\caption{Notation for Section~\ref{sec:tasks}}\label{tab:tasks-notation}
\small
\begin{tabular}{@{}ll@{\qquad}ll@{}}
\toprule
Symbol & Meaning & Symbol & Meaning\\
\midrule
$v_j,\ t_j,\ w_j$ & volume, hours per unit, loaded wage & $x_j,\ \mathcal M_j$ & mode, feasible modes\\
$c_j(x;\pi)$ & unit cost of mode $x$ & $\underline c_j$ & best alternative to $A$\\
$\pi,\ n_j(x)$ & inference price, units per unit of work & $\pi^*_j,\ v^*_j$ & crossover price, volume threshold\\
$e_j,\ L_j$ & raw AI error rate, loss per error & $\xi_j,\ \xi^{\min}_j$ & review share, oversight floor\\
$t^R_j,\ t^U_j$ & review and augmented hours per unit & $\bar q_j,\ \bar\varepsilon_j$ & required reliability, $1-\bar q_j$\\
$d^R_j$ & detection rate of review & $e^U_j$ & error rate added by augmentation\\
$c^Q_j$ & quality cost (review plus errors) & $\iota_j,\ \underline\iota_j$ & data readiness, feasibility floor\\
$\eta_e(y,y')$ & handoff cost per unit & $\eta^\times_e$ & integration premium of edge $e$\\
$q_j(y),\ Q_p$ & step and process reliability & $\Xi_p$ & process overhead\\
$W_p(\mathcal Z)$ & value of automating the set $\mathcal Z$ & $h_p$ & integration constant\\
$K_j(x),\ \delta_K$ & capital per unit of volume, replacement rate & $T_{\mathrm{par}},\ \varrho(t)$ & parallel run, ramp\\
$\varpi_p,\ b_p$ & success probability, realized share & $\mathcal I_T$ & information at data tier $T$\\
$\delta_M,\ T_{\mathrm{up}}$ & maintenance rate, model-upgrade cycle & $\Omega_j,\ C^{\mathrm{sec}}_p$ & incident exposure, security spending\\
$e^o_j,\ e^h_j,\ e^m_j$ & error, hallucination, misbehavior rates & $\mathsf a,\ c_{\mathrm{err}}$ & autonomy level, error-loss floor\\
\bottomrule
\end{tabular}
\end{table}

\subsection{The cost frontier of an activity}\label{sec:tasks-frontier}

\begin{proposition}[Mode choice and the right benchmark]\label{prop:tasks-modes}
Fix activity $j$ and its feasible set.
\begin{enumerate}
\item[(a)] The cost frontier $c^*_j(\pi):=\min_{x\in\mathcal M_j}c_j(x;\pi)$ is nondecreasing, concave and piecewise affine in $\pi$. Where the optimal mode is unique, the slope of $c^*_j$ is that mode's inference intensity.
\item[(b)] (Mode ladder.) If $\pi'<\pi$, every optimal mode at $\pi'$ is at least as inference-intensive as every optimal mode at $\pi$. As inference gets cheaper, an activity moves from $H$, $O$ or $N$ to $U$ to $A$, and never back.
\item[(c)] (Benchmark.) Automation is optimal if and only if $c_j(A;\pi)\le\underline c_j(\pi)$, and it is worth $v_j[\underline c_j-c_j(A)]^+$ a year. A business case built against in-house cost overstates this by an amount in $[0,\,v_j(c_j(H)-\underline c_j)]$; if $c_j(O)<c_j(A)<c_j(H)$, it recommends a mode that loses at least $v_j[c_j(A)-c_j(O)]$ a year.
\item[(d)] (Crossover price.) Automation is optimal if and only if $\pi\le\pi^*_j:=\min_{x\in\mathcal M_j\setminus\{A\}}[c^0_j(x)-c^0_j(A)]/[n_j-n_j(x)]$.
\item[(e)] (Volume threshold.) Let $c^{A,\mathrm{var}}_j:=c_j(A)-F^A_j/(v_jT_A)$ and suppose the alternatives have no fixed cost. Automation is optimal if and only if $\underline c_j>c^{A,\mathrm{var}}_j$ and $v_j\ge v^*_j:=F^A_j/[T_A(\underline c_j-c^{A,\mathrm{var}}_j)]$.
\end{enumerate}
\end{proposition}

\noindent\emph{Intuition.} Each mode's cost is a line in the inference price, and the frontier is their lower envelope. Cheaper inference favors the modes that use more of it, a single-crossing property, so the ladder runs from human modes through augmentation to automation; the error term of augmentation does not depend on $\pi$, so the ladder survives it, but augmentation may never be on the frontier. Part (c) is the answer to the question of what automation is worth when outsourcing is cheaper than AI: its value is $\min\{c_j(H),c_j(O)\}-c_j(A)$, not $c_j(H)-c_j(A)$. In the running example of Table~\ref{tab:tasks-example} below, matching costs \icMatchH{} in house and \icMatchO{} outsourced. Automated at readiness \icIotaHigh{} with a detection rate of \icDill{} (the conditional case of Section~\ref{sec:tasks-review}), it costs \icMatchA, so a case built against in-house cost shows a saving of \icMatchVsH{} a year on \icVolume{} invoices, while automation in fact loses \icMatchVsO{} a year against the vendor. Volume thresholds are why technical exposure overstates adoption: only 23\% of the wages paid for vision tasks exposed to AI are attractive to automate at current costs \citep{svanberg2024beyond}.

\begin{proposition}[The automation curve]\label{prop:tasks-curve}
Hold vendor prices fixed, and let $S_A(\pi)$ be the labor-equivalent spend share of the activities for which automation is optimal at $\pi$.
\begin{enumerate}
\item[(a)] The automation curve is
\[
S_A(\pi)=\frac{\sum_jE_j\mathbf 1\{\pi\le\pi^*_j\}}{\sum_jE_j},
\]
a nonincreasing step function: the spend-weighted survival function of crossover prices.
\item[(b)] (Outsourcing shifts the curve.) Adding the outsourcing option replaces $\pi^*_j$ by $\min\{\pi^*_j,[c_j(O)-c^0_j(A)]/n_j\}$, so $S_A$ falls at every $\pi$. If $\pi$ falls at the constant rate $\delta_\pi$, each crossover date is delayed by $\delta_\pi^{-1}\ln[\pi^*_j/\pi^*_j(O)]$ years, or forever if $\pi^*_j(O)\le 0$.
\item[(c)] (Value of cheaper inference.) Annual savings from AI, $\mathcal W(\pi):=\sum_jv_j[c^{\mathrm{pre}}_j-c^*_j(\pi)]$, are nonnegative, nonincreasing and convex in $\pi$, and $-\mathcal W'(\pi)=N(\pi):=\sum_jv_jn_j(x^*_j(\pi))$, the firm's inference demand, wherever optimal modes are unique.
\end{enumerate}
\end{proposition}

\noindent\emph{Intuition.} The curve inherits its shape from the distribution of crossover prices across a firm's spend: lumpy for a firm with a few large activities, smooth for one with many small ones. Outsourcing raises the bar automation must clear, so it moves each activity's crossover to a lower price and, along a falling price path, to a later date. Part (c) is Shephard's lemma: the value of a fall in the inference price is the inference the firm buys. If the price elasticity of inference demand is near one, inference spend $\pi N(\pi)$ is roughly constant as prices fall and Part I's compute share $c_c$ need not decline with $\pi$; if it is below one, spend falls with the price. The only API-data estimate, a preliminary $-1.1$ across providers of identical open models \citep{demirer2025emerging}, measures substitution between providers, and its authors report that aggregate elasticities are substantially smaller, so the near-one case is a scenario.

Outsourced activities have their own decline path, because vendors adopt AI too. Let the vendor's unit cost be $c^V_j(\pi)=\min\{c^{VH}_j,\pi n_j+c^{V0}_j\}$, where $c^{VH}_j$ is its human cost and $c^{V0}_j$ its non-inference AI cost (low, because it spreads builds over many clients), and let its price be $p_{Oj}(\pi)=(1+\upsilon_O)c^V_j(\pi)$ with markup $\upsilon_O\ge 0$. The vendor switches to AI at $\pi^V_j:=(c^{VH}_j-c^{V0}_j)/n_j$.

\begin{proposition}[Vendor AI and the in-house automation window]\label{prop:tasks-window}
Let $D_j(\pi):=c_j(O;\pi)-c_j(A;\pi)$.
\begin{enumerate}
\item[(a)] $D_j$ is continuous, nondecreasing on $[0,\pi^V_j]$ and strictly decreasing on $[\pi^V_j,\infty)$. In-house automation therefore beats outsourcing on an interval $[\pi^-_j,\pi^+_j]$ that contains $\pi^V_j$, with $\pi^+_j=[(1+\upsilon_O)c^{VH}_j+\kappa_Oc_j(H)+\zeta_{Oj}-c^0_j(A)]/n_j$. The interval is nonempty if and only if $c_j(A;\pi^V_j)\le c_j(O;\pi^V_j)$: the firm must beat its vendor's pre-AI price at the moment the vendor itself switches.
\item[(b)] (Transitory automation.) If $D_j(0)=(1+\upsilon_O)c^{V0}_j+\kappa_Oc_j(H)+\zeta_{Oj}-c^0_j(A)\ge 0$, then $\pi^-_j=0$. Otherwise $\pi^-_j=-D_j(0)/(\upsilon_On_j)$ when $\upsilon_O>0$ (the interval is empty when $\upsilon_O=0$), and below $\pi^-_j$ the AI-enabled vendor is cheaper.
\item[(c)] (Make or buy at scale.) $D_j(0)\ge 0$ if and only if $v_j\ge v^{MB}_j:=F^A_j/\{T_A[(1+\upsilon_O)c^{V0}_j+\kappa_Oc_j(H)+\zeta_{Oj}-c^{A,\mathrm{var}}_j(0)]\}$ when the bracket is positive; otherwise the vendor wins in the long run at any volume. The threshold $v^{MB}_j$ is decreasing in the coordination share $\kappa_O$ and the premium $\zeta_{Oj}$ and increasing in the build cost $F^A_j$.
\item[(d)] If $\pi$ falls at rate $\delta_\pi$, the window lasts $\delta_\pi^{-1}\ln(\pi^+_j/\pi^-_j)$ years. With modes $H$, $O$ and $A$, the in-house automated share need not be monotone in $\pi$, but the AI-executed share (in-house automation or an AI-enabled vendor) is nonincreasing in $\pi$.
\end{enumerate}
\end{proposition}

\noindent\emph{Intuition.} The incumbent races its vendor. Before the vendor adopts AI, its price embeds a human cost and falling inference prices widen the incumbent's advantage. Afterwards the two cost lines fall together, and the vendor's markup on inference shrinks while its scale advantage in fixed costs remains, so the gap closes. The decision is about timing: in-house automation pays only if the window is long enough to recover the build cost, and below the make-or-buy volume $v^{MB}_j$ the firm should wait for the vendor. The window can also be empty. In the invoice chain, review and the build alone cost \icWinCzeroA{} per matched invoice in the conditional case, more than the vendor's all-in price of \icMatchO, so in-house automation of matching beats the vendor at no inference price, and a vendor that must meet the same reliability target has no reason to switch to AI either; the window opens only where review is cheap relative to the vendor's price. Part (c) separates two forces that AI exerts on the boundary of the firm. Agents that search, negotiate and transact make markets cheaper to use, a fall in $\kappa_O$ that raises $v^{MB}_j$ and favors the vendor \citep{shahidi2025coasean}; AI also lowers the cost of acquiring knowledge in house, a fall in $F^A_j$ relative to the vendor's pooled cost that lowers $v^{MB}_j$ and favors the firm \citep{garicano2026vanishing}. Which force dominates is an empirical question. The vendor's price need not fall uniformly either: after generative AI, freelance demand for automation-prone work fell, and the remaining jobs were more complex and better paid \citep{hui2024short,demirci2025who}, so the price of the residual outsourced work of an activity can rise as its easy units are automated.

\subsection{Automation, augmentation and human review}\label{sec:tasks-review}

\begin{assumption}[Errors, stakes and review]\label{ass:tasks-review}
In mode $A$, each unit of activity $j$ is wrong with probability $e_j\in(0,1]$, independently of its stakes. An undetected error costs $L_jz$, where the stakes multiplier $z\sim\mathrm{Exp}(1)$ is observed before review (an invoice amount, an account value, a severity class). Reviewing a unit costs $w_jt^R_j$, with $t^R_j\le t_j$, and detects its error with probability $d^R_j\in(0,1]$, independently of the stakes. The firm reviews a share $\xi_j$ of units, subject to a reliability requirement, that the undetected error rate $\varepsilon_j$ not exceed $\bar\varepsilon_j:=1-\bar q_j$, and an oversight floor $\xi_j\ge\xi^{\min}_j$. The floor is a modeling assumption that represents a human-oversight requirement; the AI Act requires that high-risk systems allow effective human oversight but sets no minimum review share \citep{eu2024regulation}. The raw error rate $e_j=e_j(\iota_j)$ is continuous and strictly decreasing in data readiness $\iota_j\in[0,1]$, the share of the activity's inputs that are machine-accessible: digitized, permissioned, joinable and of known quality.
\end{assumption}

\begin{proposition}[Optimal review]\label{prop:tasks-review}
Under Assumption~\ref{ass:tasks-review}:
\begin{enumerate}
\item[(a)] It is optimal to review the highest-stakes units. With review share $\xi$, the quality cost is $c^Q_j(\xi)=w_jt^R_j\xi+e_jL_j[1-d^R_j\xi(1+\ln(1/\xi))]$, which is strictly convex, and the undetected error rate is $\varepsilon_j=e_j(1-d^R_j\xi)$. When the reliability target can be met, the optimal share is
\begin{equation}\label{eq:tasks-xi}
\xi^*_j=\max\Bigl\{\exp\Bigl(-\frac{w_jt^R_j}{d^R_je_jL_j}\Bigr),\ \frac{1-\bar\varepsilon_j/e_j}{d^R_j},\ \xi^{\min}_j\Bigr\},
\end{equation}
the largest of an economic, a reliability and a regulatory requirement.
\item[(b)] If only the economic term binds, $c^Q_j=e_jL_j[1-d^R_j\exp(-w_jt^R_j/(d^R_je_jL_j))]<\min\{e_jL_j,\,w_jt^R_j+(1-d^R_j)e_jL_j\}$.
\item[(c)] $\xi^*_j$ is nondecreasing in $e_j$, $L_j$, $\bar q_j$ and $\xi^{\min}_j$ and nonincreasing in $w_jt^R_j$; a higher detection rate raises the economic share and lowers the reliability share. The quality cost is strictly increasing in $e_j$ wherever $\xi^*_j<1$, with $dc^Q_j/de_j=L_j[1-d^R_j\xi^*_j(1+\ln(1/\xi^*_j))]$ when the economic term binds, and nonincreasing in $d^R_j$; the undetected error rate is nondecreasing in $e_j$. Review, quality cost and residual errors therefore all fall as readiness rises.
\item[(d)] (Reliability and detection thresholds.) The reliability requirement is slack if and only if $e_j\le\bar\varepsilon_j$. It can be met at some review share if and only if $d^R_j\ge 1-\bar\varepsilon_j/e_j$, that is, $e_j(1-d^R_j)\le\bar\varepsilon_j$: below this detection rate no amount of review meets the target, and automation is infeasible however cheap review and inference become. Because $e_j$ falls with readiness, this sets a readiness floor $\underline\iota_j$ at which $e_j(\underline\iota_j)(1-d^R_j)=\bar\varepsilon_j$. As $e_j\to 0$, $\xi^*_j\to\xi^{\min}_j$ (faster than any power of $e_j$ if $\xi^{\min}_j=0$) and $c^Q_j\to\xi^{\min}_jw_jt^R_j$, so an oversight floor puts a floor $\xi^{\min}_jw_jt^R_j$ under the unit cost of automation.
\end{enumerate}
\end{proposition}

\noindent\emph{Intuition.} A unit is worth reviewing when the expected loss that review averts, $d^R_je_jL_jz$, exceeds the review cost, so the economic share is the probability that stakes exceed $w_jt^R_j/(d^R_je_jL_j)$. The quality cost is then a smooth minimum of reviewing everything and reviewing nothing. Reliability targets and oversight rules can only add review, and imperfect detection makes each reviewed unit buy less reliability. Field and laboratory evidence shows that reviewers miss many AI errors: radiologists' accuracy fell from about 80\% to between 20\% and 46\%, depending on experience, when the AI's suggestion was wrong \citep{dratsch2023automation}; human-AI combinations do worse than the better of the two alone in decision tasks \citep{vaccaro2024when}; and radiologists given AI predictions underweight them \citep{agarwal2023combining}. Measured catch rates are about 0.63 for an individual inspector and 0.83 for a group on seeded ordinary errors \citep{panko1999applying}, and 0.20 to 0.46 for confident wrong AI output \citep{si2024large,dratsch2023automation}. The value of review therefore depends on its detection rate as well as its cost, and below the threshold of part (d) no amount of review meets a reliability target. This is the binding constraint in the running example. At the engine's defaults, the raw error rate at readiness \icIotaHigh{} is \icEhighMin{} to \icEhighMax{} and reviewers catch \icDR{} of errors, so no step can meet its target at any review share: the detection thresholds are \icDstarMin{} to \icDstarMax. Capture and matching, whose target as finance steps is the engine's human-parity level of \icQbarHigh, reach it at the default catch rate only above readiness \icCaptureIotaFeas{} and \icMatchIotaFeas; approval, whose activity prior asks for \icQbarAppr, never does, since it reaches at most \icMaxRelApprFull{} at full readiness, where the raw error rate is still \icEfullMin; exception handling, with a \icQbarExc{} target, becomes automatable only above readiness \icExcIotaFeas. The whole chain needs a detection rate of \icDfeas{} at readiness \icIotaHigh, above the measured range of \icDRlow{} to \icDRhigh. The process results below are therefore illustrated in a conditional case. Detection is set at the top of the measured range, \icDill, and every step's reliability target, which is the firm's policy choice rather than an empirical value, is capped at the human-parity level of \icQbarHigh: the firm asks the automated chain to be as reliable as top-performing accounts-payable staff and no more. The cap binds only on approval, which accepts a residual error of \icResidCap{} instead of \icResidAppr, equal to the human baseline error rate of \icEHuman; it raises no step's target, and if the chain could not meet human parity the cap would fall to the highest level it meets. The loss per undetected error stays at \icLmultReg{} times the labor cost, so the expected error loss rises with the residual. Every other input is at its default. All four steps stay regulated, and their builds keep the regulated multiplier of \icRegMult: accounts payable is in scope for internal control over financial reporting at a public company, and resolving three-way-match exceptions changes recorded liabilities. Regulated status is a property of the process, which the firm's error tolerance does not change. At readiness \icIotaHigh{} the review shares are \icXiMin{} to \icXiMax: the ``automation'' is almost entirely reviewed. The \icQbarHigh{} target rests on benchmark reports that put top performers' first-time error-free disbursements at about 98\% (a secondary report of APQC data) and a vendor's best customers at 99.1\% \citep{fraxion2025apbenchmarks,medius2025apaccuracy}, so it is class b with a range of \icRbarHighLow{} to \icRbarHighHigh; the loss of \icLmultReg{} times the labor cost per error is a scenario value (class d), with a range of \icLmultRegLow{} to \icLmultRegHigh{} times. Together with the detection rate and the run costs of Section~\ref{sec:tasks-ancillary}, they decide whether the chain is feasible and the sign of its value. Capping approval instead at \icQstar, the highest common target that capture, matching and approval meet at this detection rate, asks more of the automated step than of people, and moves the process results below by less than \$10{,}000 a year: integration is then worth \icAltWfull{} at readiness \icIotaHigh.

\begin{proposition}[Augment or automate: the activity readiness threshold]\label{prop:tasks-augment}
Suppose augmentation is feasible at every readiness level at a cost that does not depend on $\iota_j$, automation is feasible if and only if $\iota_j\ge\underline\iota_j$, and $c_j(U)<\min\{c_j(H),c_j(O)\}$.
\begin{enumerate}
\item[(a)] $c_j(A;\iota_j)$ is continuous and strictly decreasing in $\iota_j$.
\item[(b)] There is a threshold $\iota^*_j=\max\{\underline\iota_j,\iota^c_j\}$, where $c_j(A;\iota^c_j)=c_j(U)$ ($\iota^c_j=0$ if $c_j(A;0)\le c_j(U)$ and $\iota^c_j=\infty$ if $c_j(A;1)>c_j(U)$), such that augmentation is optimal for $\iota_j<\iota^*_j$ and automation for $\iota_j\ge\iota^*_j$.
\item[(c)] Automation beats augmentation only if $\xi^*_jw_jt^R_j\le w_jt^U_j+e^U_jL_j-\pi(n_j-n^U_j)-F^A_j/(v_jT_A)-(r+\delta_K)[K_j(A)-K_j(U)]$. In particular, if reviewing a unit takes as long as doing it with assistance ($t^R_j=t^U_j$) and $e^U_j=0$, automation requires a review share below one by the inference, build and capital gap, and an oversight floor above that level keeps the activity augmented at every readiness level.
\item[(d)] (Decision and creation tasks.) Let $t^R_j=t^U_j$ and the reliability target be attainable. If
\[
[e^U_j-(1-d^R_j)e_j]L_j\ \ge\ \pi(n_j-n^U_j)+F^A_j/(v_jT_A)+(r+\delta_K)[K_j(A)-K_j(U)],
\]
automation with optimal review costs no more than augmentation, so augmentation is never strictly optimal. Where an assisted person errs more often than fully reviewed AI by more than the inference, build and capital gap, as in decision tasks in which the AI is more accurate, blended augmentation is dominated and the relevant comparison is automation with review routing against in-house work and outsourcing. Augmentation remains the natural mode where people outperform the AI ($e^U_j<0$) or assisted work is faster than review ($t^U_j<t^R_j$), as in creation tasks.
\end{enumerate}
\end{proposition}

\noindent\emph{Intuition.} Augmentation keeps a person on every unit, so it needs little data readiness and gains little from it; field evidence on augmentation is in \citet{brynjolfsson2025generative}. Its quality cost depends on the task: inside the AI's frontier assistance improves work, and outside it people follow wrong suggestions \citep{dellacqua2026navigating}, while human-AI combinations lose in decision tasks and gain in creation tasks \citep{vaccaro2024when}. Automation pays over augmentation only when most units can skip the person, which requires readiness, and when reviewers catch enough errors. In the running example, augmented decision steps make \icEUdecision{} percentage points more errors than in-house work, and exception handling costs \icExcU{} augmented. Automated at readiness \icIotaHigh, it would cost \icExcAone{} with perfect detection and \icExcAill{} at a detection rate of \icDill, the top of the measured range; it ties augmentation at \icExcDtie{} and cannot meet its target below \icExcDinf, which excludes the default rate of \icDR. At the default rate its activity threshold is $\iota^*=\icExcIotaStarDefault$.

\subsection{Processes: handoffs, complementarity and chained errors}\label{sec:tasks-process}

\begin{assumption}[Process cost]\label{ass:tasks-process}
Fix each activity's best non-automated mode $M$, with $c_j(M)=\underline c_j$. For an automated set $\mathcal Z\subseteq\mathcal J_p$ let $y_j=A$ if $j\in\mathcal Z$ and $y_j=M$ otherwise. The annual cost of process $p$ is
\begin{equation}\label{eq:tasks-process}
C_p(\mathcal Z)=\sum_{j}v_jc_j(y_j)+\sum_{e=(i,j)\in\mathcal E_p}v_e\,\eta_e(y_i,y_j)+v_pL_p\bigl[1-Q_p(\mathcal Z)\bigr]+\Xi_p\mathbf 1\{\mathcal C_p\not\subseteq\mathcal Z\},
\end{equation}
with $Q_p(\mathcal Z)=\prod_jq_j(y_j)$. Here $\eta_e$ is the handoff cost per unit (re-keying, queueing, reconciliation, and latency valued at a rate per unit); $q_j(y)$ is the probability that the step succeeds, independently across steps, so a unit of output fails with probability $1-Q_p$, as in the O-ring model of \citet{kremer1993oring} and its task-level version in \citet{gans2026oring}; $L_p$ is the loss per failed unit beyond the local losses in $c_j$; and $\Xi_p$ is the process's fixed overhead (supervision, minimum staffing, queue management), released only when the whole critical path is automated. The value of automating $\mathcal Z$ is $W_p(\mathcal Z):=C_p(\emptyset)-C_p(\mathcal Z)$, and the integration premium of edge $e$ is $\eta^\times_e:=\eta_e(A,M)+\eta_e(M,A)-\eta_e(A,A)-\eta_e(M,M)$.
\end{assumption}

\begin{proposition}[Complementarity]\label{prop:tasks-complement}
Under Assumption~\ref{ass:tasks-process}:
\begin{enumerate}
\item[(a)] If $\eta^\times_e\ge 0$ for every edge and automation moves all step reliabilities in the same direction ($q_j(A)\ge q_j(M)$ for all $j$, or $q_j(A)\le q_j(M)$ for all $j$), then $W_p$ is supermodular: $W_p(\mathcal Z\cup\mathcal Z')+W_p(\mathcal Z\cap\mathcal Z')\ge W_p(\mathcal Z)+W_p(\mathcal Z')$.
\item[(b)] The handoff part of $W_p$ is $\sum_{j\in\mathcal Z}W^\eta_p(\{j\})+\sum_{e\subseteq\mathcal Z}v_e\eta^\times_e$, where $e\subseteq\mathcal Z$ means both ends are automated.
\item[(c)] (Monotone expansion.) Under (a), the largest optimal automated set expands as $\pi$ falls and as the set of feasible activities grows.
\end{enumerate}
\end{proposition}

\noindent The sign condition $\eta^\times_e\ge 0$ is not automatic. AI predictions raise the variance of decisions, which is costly when the decisions at the two ends of an edge interact \citep{agrawal2024artificial}; automating both ends of such an edge can then raise the handoff cost, and its premium can be negative. Supermodularity then requires a coordination mechanism (shared decision rules, joint review of the two steps), whose cost belongs in the integration build $F^A$, as do any reductions in interdependence that give up synergies.

\begin{proposition}[All-or-nothing]\label{prop:tasks-aon}
Suppose (a) of Proposition~\ref{prop:tasks-complement} holds, $\mathcal C_p=\mathcal J_p$, and $q_j(A)/q_j(M)$ is the same for all $j$. Define the integration constant
\begin{equation}\label{eq:tasks-h}
h_p:=\min_{\emptyset\ne\mathcal Z\subsetneq\mathcal J_p}\Bigl[\frac{|\mathcal Z|}{J}\sum_{e\in\mathcal E_p}v_e\eta^\times_e-\sum_{e\subseteq\mathcal Z}v_e\eta^\times_e\Bigr]\frac{J}{|\mathcal Z|\,(J-|\mathcal Z|)}.
\end{equation}
\begin{enumerate}
\item[(a)] If $\max_jW_p(\{j\})-\min_jW_p(\{j\})\le h_p+\Xi_p/(J-1)$, then $W_p(\mathcal Z)\le(|\mathcal Z|/J)\,W_p(\mathcal J_p)$ for every $\mathcal Z$. The optimal automated set is then $\emptyset$ if $W_p(\mathcal J_p)<0$ and $\mathcal J_p$ if $W_p(\mathcal J_p)>0$.
\item[(b)] If $v_e\eta^\times_e=c$ on every edge, $h_p=c/(J-1)$ for every tree (chains and assembly trees), $h_p=4c/J$ for a cycle with $J$ even ($4cJ/(J^2-1)$ for $J$ odd) and $h_p=cJ/2$ for a complete graph.
\item[(c)] (Sharpness.) For a chain with $v_e\eta^\times_e=c$, $\Xi_p=0$ and unchanged reliability, the bound $c/(J-1)$ cannot be raised: for every dispersion above it, there are stand-alone values with $W_p(\mathcal J_p)=0<W_p(\mathcal J_p\setminus\{J\})$.
\end{enumerate}
\end{proposition}

\noindent\emph{Intuition.} A lone automated step in a manual chain pays the cost of a mode boundary on each side, and the premium is recovered only when its neighbors are automated too; chained errors and the overhead released by straight-through processing add to the complementarity \citep{milgrom1990economics,topkis1998supermodularity}. When coupling is strong relative to the heterogeneity of stand-alone values, a process is automated whole or not at all. The integration constant is smallest for trees, where the end activities share a single handoff: long chains peel at their ends, so partial automation of a long chain starts with its first or last steps, and the tolerated dispersion falls like $1/(J-1)$. Rework loops, which close cycles, raise it. \citet{gans2026oring} reach the same conclusion, that adoption is discrete and bundled, from quality complementarity in a worker's allocation of time; the condition here adds handoffs, chained errors and released overhead on a process graph, and the integration constant says how much heterogeneity a given graph tolerates.

\begin{proposition}[The value of full integration]\label{prop:tasks-integration}
Under Assumption~\ref{ass:tasks-process}:
\begin{enumerate}
\item[(a)] $W_p(\mathcal J_p)=\sum_jv_j[c_j(M)-c_j(A)]+\sum_ev_e[\eta_e(M,M)-\eta_e(A,A)]+v_pL_p[Q_p(\mathcal J_p)-Q_p(\emptyset)]+\Xi_p$: unit-cost, coordination, reliability and overhead dividends.
\item[(b)] (Piecemeal cases understate integration.) With $\hat q_j:=q_j(A)/q_j(M)$ and $|\mathcal C_p|\ge 2$,
\begin{equation}\label{eq:tasks-premium}
W_p(\mathcal J_p)-\sum_jW_p(\{j\})=\sum_ev_e\eta^\times_e+v_pL_pQ_p(\emptyset)\Bigl[\prod_j\hat q_j-1-\sum_j(\hat q_j-1)\Bigr]+\Xi_p,
\end{equation}
which is nonnegative under (a) of Proposition~\ref{prop:tasks-complement}; with negative premia at interdependent edges and no coordination mechanism, piecemeal cases can instead overstate it.
\item[(c)] Full integration beats a cheaper manual chain, even one that is cheaper at every step, if and only if the unit-cost penalty $\sum_jv_j[c_j(A)-c_j(M)]$ is smaller than the coordination, reliability and overhead dividends.
\end{enumerate}
\end{proposition}

\begin{lemma}[Chained errors]\label{lem:tasks-oring}
(a) An automated chain of $J$ steps with residual errors $\varepsilon_k$ meets a process reliability target $\bar q_p$ if and only if $\sum_{k=1}^J-\ln(1-\varepsilon_k)\le-\ln\bar q_p$. With a common error $\varepsilon$ the condition is $\varepsilon\le 1-\bar q_p^{1/J}$, and $(1-\bar q_p)/J\le 1-\bar q_p^{1/J}\le-\ln(\bar q_p)/J$. If $\varepsilon_k$ is nondecreasing in the position $k$, the longest chain that meets the target is no longer than the common-error bound computed with $\varepsilon_1$. (b) $\partial Q_p/\partial q_j=Q_p/q_j$ and $\partial^2Q_p/\partial q_i\partial q_j=Q_p/(q_iq_j)>0$: reliability gains are complements, and the marginal value of reliability is highest at the least reliable step.
\end{lemma}

\noindent\emph{Intuition and example.} A ten-step agentic chain that must succeed 99\% of the time needs a per-step error rate below 0.1\%. The common-error bound is optimistic for long single-agent chains, because per-step accuracy falls as steps accumulate when models condition on their own earlier errors \citep{sinha2025illusion} and errors compound across compositional steps \citep{dziri2023faith}. Table~\ref{tab:tasks-example} sets out the running example, a four-step supplier-invoice chain whose every input is an engine default or a stated example setting (\texttt{data/examples/invoice\_chain.yaml}). The firm is a standalone business transforming in place: its integration cost, readiness and change capacity are its own. At readiness \icIotaHigh, its cheapest manual configuration eliminates capture, which e-invoicing makes possible at that readiness, outsources matching and augments exception handling and approval, at \icCmanHigh{} a year including handoffs, against \icEinhouse{} in house. In the conditional case (detection \icDill{} and approval capped at \icQcap), the unit-cost dividend of full automation is \icDivUnit, because only exception handling is cheaper automated; the coordination dividend is \icDivCoord, the released overhead \icDivOver{} and the reliability dividend \icDivRel, as process reliability falls from \icQM{} to \icQA. Integration is worth \icWfullHigh{} a year. Automated alone, every step loses money (capture \icSingleCapture, matching \icSingleMatch, exception handling \icSingleExc{} and approval \icSingleAppr), but exception handling and approval together gain \icPairEA{} a year, the best set at readiness \icIotaHigh. The pair is a small instance of complementarity: it is worth more than its two steps alone and more than the whole chain, and it is the first set to pay, above readiness \icBestFirstIota, before the whole chain (above \icIotaStarN). Its sign at \icIotaHigh{} is fragile: with approval capped at \icQstar{} instead, the pair loses \icAltPairLoss{} a year and first pays above \icAltFirstIota. The whole chain is not all-or-nothing, because capture is better eliminated and matching better outsourced, a dispersion of stand-alone values far above the integration constant. At the firm level, firm performance is positively associated with the breadth of AI integration across business functions \citep{bonney2026microstructure}, an association rather than an effect, consistent with the integration premium.

\begin{table}[t]
\centering
\caption{Running example: a four-step supplier-invoice chain, \icVolume{} invoices a year (engine defaults)}\label{tab:tasks-example}
\footnotesize
\setlength{\tabcolsep}{4pt}
\begin{tabular}{@{}lccccccccc@{}}
\toprule
Step & $t_j$ (h) & $c_j(H)$ & $c_j(M)$, mode & $\pi n_j$ & $e_j(\icIotaLow)$ / $e_j(\icIotaHigh)$ & $\bar q_j$ & $d^{R*}_j$ & $c_j(A)$, cond. & $F^A_j$ (\$000)\\
\midrule
\icRowCapture\\
\icRowMatch\\
\icRowExceptions\\
\icRowApproval\\
\bottomrule
\end{tabular}

\smallskip
\begin{minipage}{0.97\linewidth}\footnotesize
Notes: wage \icWage{} an hour (OEWS median \icWageAnnual{} for occupation 43-3031, loaded); dollar unit costs per invoice; $c_j(M)$, $\pi n_j$, $\bar q_j$, the detection threshold $d^{R*}_j$ and $c_j(A)$ at readiness \icIotaHigh. Vendor price \icRhoO{} of in-house cost, $\kappa_O=\icKappaO$, $\zeta_O=\icZetaO\,w_jt_j$; outsourcing infeasible for exceptions and approval and elimination feasible for capture above readiness \icIotaN{} (engine priors); $t^U_j=t_j/(1+\icBetaU)$; review takes $\icRhoRev\,t_j$ at \icWagePremium{} times the wage and detects errors with probability $d^R=\icDR$; the conditional case (cond.) has $d^R=\icDill$ and targets of at most \icQcap; $L_j$ is \icLmult{} times $w_jt_j$; $e_j(\iota)$ combines $e^1+(e^0-e^1)(1-\iota)$ ($e^0=\icEzero$, $e^1=\icEone$) with the capability error at a 50\% horizon of \icHzero{} hours, and capture runs through a user interface below readiness \icApiReadiness; builds by the engine's rule, with its regulated multiplier of \icRegMult{} for all four steps (a financial-reporting process), amortized at the capital recovery factor ($r=\icR$, $T_A=\icTA$). Handoffs per invoice and edge: $\eta(M,M)={}$\icEtaMM, $\eta(A,M)=\eta(M,A)={}$\icEtaAM, $\eta(A,A)={}$\icEtaAA$(1-\iota)$. Process failure loss $L_p={}$\icLp, $q_j(M)=\icQMstep$, overhead $\Xi_p=\icXiPshare$ of in-house spend. Inputs and sources: \texttt{data/examples/invoice\_chain.yaml}; values: \texttt{code/examples/invoice\_chain.py}.
\end{minipage}
\end{table}

\subsection{Data readiness}\label{sec:tasks-readiness}

\begin{lemma}[Straight-through processing]\label{lem:tasks-stp}
Let each critical-path step $j$ review the units whose stakes multiplier $z_j\sim\mathrm{Exp}(1)$ exceeds $\ln(1/\xi_j)$. If the $z_j$ are associated (positively dependent), the straight-through rate $\mathrm{STP}_p$, the share of units no person touches, satisfies $\prod_j(1-\xi_j)\le\mathrm{STP}_p\le\min_j(1-\xi_j)$, with equality on the left when stakes are independent across steps and on the right when they are common (an invoice's amount).
\end{lemma}

\noindent Process readiness thus aggregates step readiness, through $\xi_j(\iota_j)$, as a product or as a weakest link, depending on how stakes are correlated along the process.

\begin{proposition}[The readiness threshold]\label{prop:tasks-readiness}
Let $\mathcal C_p=\mathcal J_p$; let readiness $\iota$ be common to the process's activities, with $e_j(\iota)$ strictly decreasing, $\eta_e(A,A;\iota)$ continuous and nonincreasing, $\xi^*_j<1$ and $q_j(A;\iota)=1-e_j(\iota)(1-\xi^*_j(\iota))$; let $\underline\iota_p:=\max_j\underline\iota_j$; and let the condition of Proposition~\ref{prop:tasks-aon}(a) hold at every $\iota$.
\begin{enumerate}
\item[(a)] $W_p(\mathcal J_p;\iota)$ is continuous and strictly increasing in $\iota$.
\item[(b)] Let $\iota^*_p:=\min\{\iota\ge\underline\iota_p:W_p(\mathcal J_p;\iota)\ge 0\}$ ($\infty$ if the set is empty). Full automation is optimal if and only if $\iota\ge\iota^*_p$. Wherever $W_p(\mathcal J_p;\iota)<0$, which includes every $\iota<\iota^*_p$ when $\iota^*_p>\underline\iota_p$, no activity is automated, and each runs in its best non-automated mode: augmentation, where it beats in-house work and outsourcing.
\item[(c)] (Weakest link.) With activity-specific readiness, full automation is feasible if and only if $\min_j(\iota_j-\underline\iota_j)\ge 0$, and the set of readiness vectors at which it is optimal is an upper set. If $q_j(A;\iota_j)=q^A(\iota_j)$ with $q^A$ increasing and concave, the marginal value of readiness for process reliability, $\partial Q_p/\partial\iota_j=Q_pq^{A\prime}(\iota_j)/q^A(\iota_j)$, is largest at the least-ready step.
\item[(d)] (Big push.) Let $\underline\iota_p<\iota^*_p<\infty$. If readiness can be raised from $\iota^0<\iota^*_p$ at an increasing cost $I_p(\iota)-I_p(\iota^0)$, no level in $(\iota^0,\iota^*_p)$ is optimal: the firm should invest to at least $\iota^*_p$ or not at all, and should invest if and only if $\max_{\iota\ge\iota^*_p}\{W_p(\mathcal J_p;\iota)/r-[I_p(\iota)-I_p(\iota^0)]\}>0$.
\end{enumerate}
\end{proposition}

\noindent\emph{Intuition and example.} Readiness lowers review, residual errors and handoff costs, so the value of full automation rises with it, but under all-or-nothing nothing is gained until it crosses zero. Data integration is therefore a threshold investment, and partial integration buys nothing. In the running example at the default detection rate no readiness level makes the chain automatable. In the conditional case the chain becomes feasible at readiness \icIotaChainFeas, approval's readiness floor. Below \icIotaN{} capture's best alternative is outsourcing, and full automation loses money, \icWbelowN{} a year just below that level. At \icIotaN{} e-invoicing makes elimination of capture possible, so the benchmark improves and the value of full automation falls by \icWdropN, to \icWaboveN; the proposition holds each activity's best alternative fixed, and here the alternative itself improves with readiness. Above \icIotaN, full automation is worth \icWfullHigh{} at \icIotaHigh{} and \icWninetyfive{} at 0.95, and turns positive above \icIotaStarN. A firm at \icIotaLow{} should either fund integration to that level or leave the chain manual; once integration and run costs are counted, it is the second (Sections~\ref{sec:tasks-transition} and \ref{sec:tasks-ancillary}). Most firms appear to sit below such thresholds: in November 2025 to January 2026, 18\% of US firms used AI in at least one business function, 57\% of adopters in three or fewer functions, and 66\% of firms used AI for task augmentation \citep{bonney2026microstructure}. Because integration costs come first and benefits later, the investment shows up as a J-curve in measured performance \citep{brynjolfsson2021productivity}; Section~\ref{sec:tasks-transition} quantifies it in score terms.

\subsection{Capital per unit of work}\label{sec:tasks-capital}

Each mode can need capital per unit of annual volume, $K_j(x)$: owned computing capacity, devices and robotics, capitalized integration builds and internal-use software. It is installed with volume and replaced at rate $\delta_K$, which is high for accelerators.

\begin{proposition}[Capital for labor at the activity level]\label{prop:tasks-capital}
Let activity $j$ switch to a mode that lowers its operating unit cost by $\Delta c_j>0$ and raises its capital per unit of annual volume by $\Delta K_j>0$.
\begin{enumerate}
\item[(a)] (Value.) For any nondecreasing volume path with $\int_0^\infty\exp(-rt)v_j(t)\,dt<\infty$, the NPV of switching, including the retrofit of existing capacity, is $[\Delta c_j-(r+\delta_K)\Delta K_j]\int_0^\infty\exp(-rt)v_j(t)\,dt$. The switch creates value if and only if $\Delta c_j>(r+\delta_K)\Delta K_j$, whatever the growth path; the NPV is $v_j[\Delta c_j-(r+\delta_K)\Delta K_j]/r$ at constant volume and $v_j(0)[\Delta c_j-(r+\delta_K)\Delta K_j]/(r-g)$ at constant growth $g<r$. Mode choice should therefore use the full unit cost $c_j(x)+(r+\delta_K)K_j(x)$.
\item[(b)] (Score.) In a steady state with volume growing at rate $g$, the switch changes the FCF margin by $v_j[\Delta c_j-(g+\delta_K)\Delta K_j]/S$ and the GAAP operating margin by $v_j[\Delta c_j-\delta_K\Delta K_j]/S$. With volume proportional to revenue, $\Delta a=-v_j\Delta c_j/S$ and $\Delta k=v_j\Delta K_j/S$, and to first order $\Delta R_X^*=-\Delta a-(g^*+\delta_K)\Delta k$: Corollary~\ref{cor:substitution} with depreciation (Remark~\ref{rem:assumptions}(c)), at the activity level.
\item[(c)] (Ranking.) GAAP charges capital at $\delta_K$, the cash Rule of 40 at $g+\delta_K$ and value at $r+\delta_K$. If $g<r$, every value-creating switch raises the cash score, but switches with $(g+\delta_K)\Delta K_j<\Delta c_j<(r+\delta_K)\Delta K_j$ raise it while destroying value. If $g>r$, switches with $(r+\delta_K)\Delta K_j<\Delta c_j<(g+\delta_K)\Delta K_j$ create value but lower the cash score. Between two modes of equal full cost, the steady-state cash score prefers the more capital-intensive one if and only if $g<r$, and the GAAP margin always does.
\end{enumerate}
\end{proposition}

\noindent\emph{Intuition.} Value charges the opportunity cost of the whole capital stock, including the retrofit of today's capacity. A steady-state margin charges only the flow of capital spending: replacement, plus the capital for new volume, which is $g\Delta K_j$ per unit. The cash score sees the retrofit only once, as a dip in the adoption year (Proposition~\ref{prop:tasks-scorepath}). This also reconciles Corollary~\ref{cor:substitution} with value: Part I's hurdle $g^*$ compares steady states, and adding the retrofit raises it to $r$. For example, owning the inference capacity of the invoice chain instead of renting it saves \icCapDc{} an invoice (owned capacity costs \icOwnedOpex{} of the rented price to run) but needs \icCapDK{} of capital per invoice of annual capacity (\$\icCapPerDollar{} per annual dollar of inference), depreciating at \icDeltaKcompute{} a year (engine defaults). At $r=\icR$ and the engine's growth of \icCapG, it destroys \icCapValue{} of value per invoice a year, lowers the cash margin by \icCapCash{} and the GAAP margin by \icCapGaap{} per invoice: with growth above the discount rate, the cash score penalizes the capital more than value does.

\subsection{Transitions: parallel running, execution risk and pilots}\label{sec:tasks-transition}

\begin{assumption}[Transition]\label{ass:tasks-transition}
At date 0, process $p$ moves from its legacy mode, with annual cash cost $C_0$, to a new mode that saves $W$ a year at full realization (net of replacement capital spending, and excluding the amortized build, which is counted in $F$). The realized share of the saving is $b\in[0,1]$.
\begin{enumerate}
\item[(i)] At date 0 the firm pays the build and integration cost $F$ (including readiness investment), capital spending $I=\sum_jv_j\Delta K_j$ and a productivity dip $C^J$.
\item[(ii)] During a parallel run on $[0,T_{\mathrm{par}})$ the new mode carries the share $\varrho(t)=\hat\varrho(t/T_{\mathrm{par}})$ of volume, with $\hat\varrho:[0,1]\to[0,1]$ nondecreasing and of mean $\bar\varrho$, and saves $bW\varrho(t)$. Legacy capacity is kept for rollback at a duplication premium $\omega_{\mathrm{par}}C_0$, reconciliation costs $\nu_{\mathrm{rec}}E_p$ ($E_p$ is the process's in-house labor-equivalent spend), and duplicate licenses cost $C^{\mathrm{dup}}$, for a parallel premium $C^{\mathrm{par}}:=\omega_{\mathrm{par}}C_0+\nu_{\mathrm{rec}}E_p+C^{\mathrm{dup}}$ a year.
\item[(iii)] At cutover, date $T_{\mathrm{par}}$, the firm pays stranded cash costs $C^{\mathrm{str}}:=(1-u_d)T_{\mathrm{sev}}C^{\mathrm{disp}}+C^{\mathrm{brk}}$: severance of $T_{\mathrm{sev}}$ years of wages on the displaced wage bill $C^{\mathrm{disp}}$, except the share $u_d$ that is redeployed, and contract break fees $C^{\mathrm{brk}}$. It writes off undepreciated legacy assets $C^{\mathrm{wo}}$, a non-cash charge.
\item[(iv)] With probability $\varpi_p$ the new mode works, that is, the attempt creates some value, and then $b$ has mean $\bar b_p\in(0,1]$; otherwise $b=0$. A failure is detected during the parallel run with probability $1-\exp(-T_{\mathrm{par}}/T_{\mathrm{det}})$, and the firm then abandons at $T_{\mathrm{par}}$ without stranded costs. An undetected failure leads to cutover and then to rollback at a loss $L_{\mathrm{roll}}$ (restoring legacy capacity and running in parallel again), valued at $T_{\mathrm{par}}$.
\item[(v)] Savings accrue until the horizon $T_h$; the discount rate is $r$, and $A_r(T):=[1-\exp(-rT)]/r$.
\end{enumerate}
\end{assumption}

The success probability depends on observable drivers. We specify it at the activity level and aggregate along the critical path,
\begin{equation}\label{eq:tasks-success}
\varpi_p=\varpi^0_p\,q(L/\Phi),\qquad\varpi^0_p=\varpi^c_p\prod_{j\in\mathcal C_p}\varpi_{pj},\qquad\operatorname{logit}\varpi_{pj}=\mathsf b_0+\mathsf b_\iota\iota_j+\mathsf b_{\mathcal H}\ln(\mathcal H/\bar d_j)-\mathsf b_R\mathbf 1\{j\text{ regulated}\}.
\end{equation}
Here $\varpi_{pj}$, the probability that the automation of activity $j$ works, depends on readiness, on the capability margin of effective AI capability $\mathcal H$ over the activity's reliability-adjusted difficulty $\bar d_j$, and on the regulatory class. In the notation of Assumption~\ref{ass:ext-capability}, $\mathcal H=H(t-L_W\mathrm{ph}_j)$ and $\bar d_j=d_j\chi(\bar q^A_j)/\kappa_{\mathrm{ctx},j}$, so that $\mathcal H\ge\bar d_j$ is that assumption's feasibility test. Two forms of the reliability adjustment $\chi$ fit the horizon data. The logistic form of that assumption puts the 80\%, 95\% and 99\% horizons 2.32, 4.93 and 7.70 doublings below the 50\% horizon; a constant hazard per unit of human task time \citep{ord2025is}, under which a task of length $d$ succeeds with probability $2^{-d/\mathcal H}$ and $\chi(q)=1/\log_2(1/q)$, puts them 1.64, 3.76 and 6.11 doublings below, and the engine carries both. Horizons are also domain specific: on tasks that operate graphical computer interfaces they are 40 to 100 times shorter than on software tasks, at similar doubling rates \citep{kwa2025how}, a lag of 5.3 to 6.6 doublings, so $\mathcal H$ is indexed by domain. The factor $\varpi^c_p$ collects process-level drivers, with a logit that falls with the exception share and the affected headcount and rises with team experience and executive sponsorship; and $q(L/\Phi)\in(0,1]$ is the congestion factor of Section~\ref{sec:ext-bandwidth}, equal to one without congestion. All coefficients $\mathsf b$ are nonnegative and, absent sources, documented assumptions.

\begin{lemma}[Success along the critical path]\label{lem:tasks-success}
(a) $\varpi_p\le\min_{j\in\mathcal C_p}\varpi_{pj}$, and $-\ln\varpi^0_p$ is a sum of hazards $-\ln\varpi_{pj}$ and $-\ln\varpi^c_p$, so each added activity lowers the success probability proportionally. (b) $\partial\ln\varpi_p/\partial z_j=1-\varpi_{pj}$ for the logistic index $z_j=\operatorname{logit}\varpi_{pj}$: raising readiness or the capability margin pays most at the activity least likely to succeed, and as the other margins grow, $\varpi_p$ tends to $q(L/\Phi)\,\varpi^c_p\,\varpi_{pj_0}$ for the weakest activity $j_0$.
\end{lemma}

\begin{proposition}[Risk-adjusted transition value]\label{prop:tasks-transition}
Under Assumption~\ref{ass:tasks-transition}, let $p_u:=\exp(-T_{\mathrm{par}}/T_{\mathrm{det}})$ and $\mathcal S_r:=\int_0^{T_{\mathrm{par}}}\exp(-rt)\varrho(t)\,dt+[\exp(-rT_{\mathrm{par}})-\exp(-rT_h)]/r$.
\begin{enumerate}
\item[(a)] The expected NPV of the transition is
\begin{multline}\label{eq:tasks-npv}
\mathrm{NPV}_p=-(F+I+C^J)-C^{\mathrm{par}}A_r(T_{\mathrm{par}})+\varpi_p\bar b_pW\mathcal S_r\\
-\exp(-rT_{\mathrm{par}})\bigl\{[\varpi_p+(1-\varpi_p)p_u]C^{\mathrm{str}}+(1-\varpi_p)p_uL_{\mathrm{roll}}\bigr\}.
\end{multline}
\item[(b)] (Transition breakpoint.) The transition pays if and only if $\mathrm{NPV}_p>0$: the risk-adjusted saving times the discounted remaining horizon, $\varpi_p\bar b_pW\mathcal S_r$, must exceed $F+I+C^J$ plus the present value of the parallel-run, stranded and rollback costs. Without discounting, the minimum horizon is $T_h^*=(1-\bar\varrho)T_{\mathrm{par}}+\{F+I+C^J+C^{\mathrm{par}}T_{\mathrm{par}}+[\varpi_p+(1-\varpi_p)p_u]C^{\mathrm{str}}+(1-\varpi_p)p_uL_{\mathrm{roll}}\}/(\varpi_p\bar b_pW)$.
\item[(c)] (Payback.) If success is certain and complete, the undiscounted payback date, after which cumulative cash flow stays nonnegative, is $t_{\mathrm{pay}}=(F+I+C^J+C^{\mathrm{str}})/W+T_{\mathrm{par}}(1-\bar\varrho+C^{\mathrm{par}}/W)$ when $t_{\mathrm{pay}}\ge T_{\mathrm{par}}$: the parallel run extends payback by $T_{\mathrm{par}}(1-\bar\varrho+C^{\mathrm{par}}/W)$.
\item[(d)] (Optimal parallel run.) Without discounting, the NPV is maximized at
\begin{equation}\label{eq:tasks-tpar}
T^*_{\mathrm{par}}=T_{\mathrm{det}}\ln\Bigl[\frac{(1-\varpi_p)(C^{\mathrm{str}}+L_{\mathrm{roll}})}{T_{\mathrm{det}}\,\hat c}\Bigr]^+,\qquad\hat c:=C^{\mathrm{par}}+\varpi_p\bar b_pW(1-\bar\varrho),
\end{equation}
where $\hat c$ is the cost of a year of parallel running, duplication plus forgone savings. At an interior optimum the residual probability of an undetected failure is $(1-\varpi_p)\exp(-T^*_{\mathrm{par}}/T_{\mathrm{det}})=T_{\mathrm{det}}\hat c/(C^{\mathrm{str}}+L_{\mathrm{roll}})$. $T^*_{\mathrm{par}}$ rises with the failure probability, the rollback loss and the stranded costs, and falls with the duplication premium, the reconciliation cost and the success probability.
\end{enumerate}
\end{proposition}

\noindent\emph{Intuition.} A parallel run is insurance bought by the year: it costs the duplication premium plus the savings forgone on the volume not yet migrated, and it pays by catching failures while rollback is still cheap. The optimum stops when a further year costs more than the expected rollback loss it averts, which sets the residual risk independently of the prior failure probability. Large IT projects show fat-tailed overruns \citep{flyvbjerg2011why}, which is why the risk terms are explicit. Take the invoice chain moving from readiness \icIotaLow{} to \icIotaHigh{} in the conditional case of Section~\ref{sec:tasks-review}, since at the default detection rate it cannot be automated at all; every other input is the engine default named in \texttt{data/examples/invoice\_chain.yaml}. Before run costs the chain saves \icWbeforeRun{} a year; the run costs of Section~\ref{sec:tasks-ancillary}, \icRunTotal{} a year, make the saving $W={}$\icW. The one-time cost is $F={}$\icF: \icFbuild{} of builds by the engine's rule, inside the range of implementation costs of one AI system for one workflow (\$165{,}000 to \$1.77 million, \citealp{svanberg2024beyond}); \icFint{} of integration, the engine's default of \icIntCost{} per system for \icNsys{} systems, convex in readiness, an assumption under review; and \icFother{} of redesign, legal review, evaluation data, change management, training and control redesign. The other inputs are $C^J={}$\icCJ{} (a \icDipShare{} productivity dip for a year on \icCzero{} of legacy cost), $C^{\mathrm{par}}={}$\icCpar{} a year (\icOmegaPar{} duplication of the legacy cost, reconciliation of \icNuRec{} of in-house spend and \icCdup{} of duplicate licenses), $C^{\mathrm{str}}={}$\icCstr{} (severance, break fees and redeployment; the survivor cost of Proposition~\ref{prop:labor-survivor}(a) is zero at the engine's defaults, because induced quits do not rise with the share displaced), $\varpi_p=\icVarpi$ from the engine's success logits (calibrated so that a reference five-step process creates some value in 70\% of attempts, Part~I's $\varpi$, which is the share of attempts that create some value and not a success rate; with $\bar b=0.67$ the expected yield of a reference attempt is $\varpi\bar b\approx0.47$, Part~I's target), $\bar b_p=\icBbar$, $T_{\mathrm{det}}=\icTdet$ years, $L_{\mathrm{roll}}={}$\icLroll{} (a year of legacy cost), $T_h=\icTh$ years (the useful life of internal-use software, three to five years in ASC 350-40 practice) and $r=\icR$. The parallel run is \icTpar{} years, the engine's floor for a financial-reporting process (the unconstrained optimum is \icTparStar{} years), and leaves a \icResid{} risk of an undetected failure. The announcement adds a per-wave charge $O_s={}$\icOs{} (Proposition~\ref{prop:labor-survivor}(b): \icExcessQuits{} more quits among the \icNsv{} FTE who stay, each replaced at \icCrep{} years of wages), which enters the wave problem of Proposition~\ref{prop:ext-waves} like the overhead $O_w$, not the NPV that follows. The risk-adjusted NPV is \icNPV, and \icNPVcertain{} if success were certain: with a negative annual saving, no one-time cost (the break-even is \icFbreakeven), horizon (\icNPVunlimited{} with savings forever) or success probability (\icNPVvarpiOne{} at $\varpi_p=1$) makes the transition pay. At horizons of three and five years it is \icNPVthThree{} and \icNPVthFive. The horizon of \icTh{} years is the useful life of a single build, while the engine's planning horizon of \icThEngine{} years covers a program that includes rebuilds, here in years \icRebuilds{} at the engine's declining build cost; over the program the NPV is \icNPVprogram. The integration cost moves the NPV but not its sign: at \icIntLow, \icIntCost{} and \icIntHigh{} per system the NPV is \icNPVaLow, \icNPVaDef{} and \icNPVaHigh{} for a firm starting at readiness \icIotaLow, and \icNPVb{} for a firm whose data are already integrated. What decides feasibility and the sign are the detection rate, the \icQbarHigh{} target and the loss of \icLmultReg{} times the labor cost per error, the last two scenario values (class d) of the engine's finance rule, and the run costs: at the default target, even perfect detection makes full automation worth \icDetWdone{} a year before run costs and \icDetWdoneRun{} after them. These are the values of a standalone firm transforming in place; Proposition~\ref{prop:acq-owner} and its example show that the transition turns positive only for an owner that brings several levers at once, such as a higher success probability, a platform shared by many businesses, reuse of most of the build and integration, or a much longer horizon. A realized share of \icBbar{} suits an automation program; for programs led by augmentation the evidence argues for a much lower $\bar b_p$, since two years after ChatGPT effects on earnings and hours above 2\% are ruled out in Danish administrative data \citep{humlum2025still}.

\begin{proposition}[The score path through a transition]\label{prop:tasks-scorepath}
Hold revenue $S$ and growth fixed, so that $\Delta R_X=\Delta m$, and consider the success path with realized share $b$.
\begin{enumerate}
\item[(a)] (Cash.) $\Delta m(t)=[bW\varrho(t)-C^{\mathrm{par}}]/S$ on $[0,T_{\mathrm{par}})$ and $bW/S$ afterwards, with one-time charges $(F+I+C^J)/S$ at date 0 and $C^{\mathrm{str}}/S$ at cutover. The flow is below its baseline exactly while $\varrho(t)<C^{\mathrm{par}}/(bW)$: for a linear ramp, during the first $T_{\mathrm{par}}C^{\mathrm{par}}/(bW)$ years, or the whole run if $C^{\mathrm{par}}\ge bW$. With $b=1$ the cumulative shortfall is recovered at the payback date $t_{\mathrm{pay}}$ of Proposition~\ref{prop:tasks-transition}(c).
\item[(b)] (GAAP.) Capital spending is depreciated rather than charged at date 0. Of the build, only the part that ASC 350-40 as amended by ASU 2025-06 allows to be capitalized (not data conversion or training, and not work done while significant uncertainty about novel functions is unresolved) is amortized, from the date the software is ready for its intended use and over its book life; the rest is expensed as incurred. Severance and the write-off of legacy assets are recognized when the standards require (ASC 420 or 712 for severance, accelerated depreciation from the commitment date for the write-off), not necessarily at cutover. Proposition~\ref{prop:acct-path} gives the resulting path. In the simple convention, the whole build is capitalized and amortized at $F/(T_AS)$ a year from date 0 and $(C^{\mathrm{str}}+C^{\mathrm{wo}})/S$ is charged at cutover; the GAAP margin then reaches its new steady state at $\max\{T_{\mathrm{par}},T_A\}$.
\item[(c)] (First year.) With annual periods, $T_{\mathrm{par}}\le 1\le T_A$ and a linear ramp, the first-year changes in the score are $\Delta R^{\mathrm{cash}}_1=-[F+I+C^J+C^{\mathrm{str}}+C^{\mathrm{par}}T_{\mathrm{par}}-bW(1-T_{\mathrm{par}}/2)]/S$ and, in the simple convention, $\Delta R^{\mathrm{GAAP}}_1=\Delta R^{\mathrm{cash}}_1+[F+I-F/T_A-C^{\mathrm{wo}}]/S$, the special case $F^{\mathrm{cap}}=F$, $t_{\mathrm{go}}=0$, $T_B=T_A$ of Proposition~\ref{prop:acct-path}(c).
\end{enumerate}
\end{proposition}

\noindent\emph{Intuition.} The dip has a depth, set by the one-time charges and the parallel premium, and a duration, set by the ramp and the payback date. Cash and GAAP place it differently: cash front-loads capital spending, while GAAP spreads the capitalized build and recognizes the restructuring charges when the standards require (Section~\ref{sec:accounting}). For the invoice firm, with \icRevenue{} of revenue, the transition (with $T_{\mathrm{par}}=0.5$ and $b=1$) lowers the cash Rule-of-40 score by \icSPcashOne{} points in the first year and by \icSPcashLater{} points a year thereafter, since its savings net of run costs are negative; under the simple convention of part (c) the GAAP score falls by \icSPgaapOne{} points in the first year and by \icSPgaapLater{} points a year while the build is amortized.

\begin{proposition}[The pilot as a stage gate]\label{prop:tasks-pilot}
A pilot costs $F_{\mathrm{pil}}$, delays the decision by $T_{\mathrm{pil}}$ and yields a signal about the transition's NPV $Y$, which equals $Y_s>0$ on success and $Y_f<0$ on failure. Let $d:=\exp(-rT_{\mathrm{pil}})$.
\begin{enumerate}
\item[(a)] (Gate rule.) After the pilot, the firm should scale if and only if $\mathbb E[Y\mid\text{signal}]\ge 0$; with a binary outcome, if and only if the posterior success probability is at least $\varpi^{\mathrm{gate}}:=|Y_f|/(Y_s+|Y_f|)$.
\item[(b)] (Value.) The pilot is worth $\mathrm{VP}=d\,\mathbb E[(\mathbb E[Y\mid\text{signal}])^+]-F_{\mathrm{pil}}-(\mathbb E Y)^+$, a value of information that is bounded as in Proposition~\ref{prop:tasks-granularity}(b).
\item[(c)] (When to pilot.) If the pilot reveals the outcome, $\mathrm{VP}=d\varpi_pY_s-F_{\mathrm{pil}}-[\varpi_pY_s+(1-\varpi_p)Y_f]^+$, and piloting pays if and only if
\begin{equation}\label{eq:tasks-pilotband}
\frac{F_{\mathrm{pil}}}{dY_s}<\varpi_p<\frac{|Y_f|-F_{\mathrm{pil}}}{|Y_f|+(1-d)Y_s}.
\end{equation}
Without delay, $\mathrm{VP}=\min\{\varpi_pY_s,(1-\varpi_p)|Y_f|\}-F_{\mathrm{pil}}$. Above the band the firm should scale without a pilot; below it, it should do neither.
\end{enumerate}
\end{proposition}

\noindent\emph{Intuition.} A pilot is an option to abandon \citep{mcdonald1986value,dixit1994investment}, valuable only when the decision is genuinely in doubt, as with staged research projects \citep{roberts1981funding}. Two recent arguments reach staged adoption by other routes: a real-options model in which pilots build firm-specific capability while the frontier moves (\citealp{tewari2026pilot}, a preprint), and the result that adopting a technology whose risks are learned over time should be gradual \citep{acemoglu2024regulating}. For the invoice transition, $Y_s={}$\icYs{} and $Y_f={}$\icYf: the transition loses money even when it succeeds, the premise $Y_s>0$ fails, and a pilot, which would cost \icFpil{} (the engine's \icPilotShare{} of the builds plus \icPilotFixed), is worth minus its cost at every success probability. The firm should do neither, and the same holds for a firm whose data are already integrated ($Y_s={}$\icYsb). The proposition bites only where success pays.

\subsection{Run costs, security and scope}\label{sec:tasks-ancillary}

Automated and augmented work carries costs beyond inference, review and the build: keeping agents and pipelines working, revalidating them when models change, managing the coupling between integrations, securing them, and getting deployments through security review.

\begin{assumption}[Run, security and review costs]\label{ass:tasks-ancillary}
In modes $A$ and $U$, activity $j$ also bears:
\begin{enumerate}
\item[(i)] maintenance of $\delta_MF^x_j$ a year (prompts and agents, evaluation suites, monitoring, data pipelines, drift correction), and a revalidation cost $C^{\mathrm{reval}}_j$ at each model upgrade or vendor deprecation, every $T_{\mathrm{up}}$ years, with annual equivalent $C^{\mathrm{reval}}_j/T_{\mathrm{up}}$ (with discounting, $1/T_{\mathrm{up}}$ becomes $r/[\exp(rT_{\mathrm{up}})-1]$);
\item[(ii)] technical debt: the process pays $c^{\mathrm{int}}$ a year per integration used by automated activities and $c^{\mathrm{cpl}}N_iN_j$ for each edge with both ends automated, where $N_j$ is the number of integrations activity $j$ uses;
\item[(iii)] incidents: in mode $x$ the activity has an exposure $\Omega_j(x)\ge 0$ that is nondecreasing in the sensitivity of its data, its write and transaction rights, the number of systems it can reach and the share of untrusted external inputs (instructions injected into documents and web pages succeed against 11\% to 48\% of attempts on undefended tool-using agents, \citealp{greshake2023not,debenedetti2024agentdojo}), with $\Omega_j(U)\le\Omega_j(A)$ because a person approves actions under augmentation. Incidents arrive at the rate $h^{\mathrm{inc}}_j=\Omega_j(x)(h^f+h^uv_j)$ a year, a part per process and a part per unit, so the expected annual incident loss without security spending is $C^{\mathrm{inc}}_j:=h^{\mathrm{inc}}_jL^{\mathrm{inc}}_j$, where the loss per incident $L^{\mathrm{inc}}_j$ includes the loss of firm value beyond out-of-pocket costs \citep{kamiya2021risk,florackis2023cybersecurity}. Security spending $C^{\mathrm{sec}}_p$ on the process multiplies every hazard by a residual factor $\mathsf g(C^{\mathrm{sec}}_p)$, with $\mathsf g(0)=1$ and $\mathsf g$ positive, decreasing and convex;
\item[(iv)] deployments that need security and compliance review arrive as a Poisson stream of rate $\nu$ and are reviewed first come, first served, with exponential review times, at capacity $\Phi_{\mathrm{sec}}>\nu$ reviews a year.
\end{enumerate}
Firm-level compliance costs (security certifications, an AI management system, risk-management frameworks, insurance) are fixed costs in $a$. Migrating each system that uses public-key cryptography to post-quantum standards, whose quantum-vulnerable algorithms are deprecated after 2030 and disallowed after 2035 \citep{moody2026transition}, is a dated one-time cost per system, of the per-process kind in Proposition~\ref{prop:tasks-ancscale}.
\end{assumption}

\begin{proposition}[Full cost and the breakpoints]\label{prop:tasks-fullcost}
Under Assumption~\ref{ass:tasks-ancillary}:
\begin{enumerate}
\item[(a)] For an activity that forms its own process, the full unit cost of automation is
\begin{equation}\label{eq:tasks-fullcost}
c_j(A)=\pi n_j+c^Q_j+(r+\delta_K)K_j(A)+\frac{F^A_j/T_A+\delta_MF^A_j+C^{\mathrm{reval}}_j/T_{\mathrm{up}}+c^{\mathrm{int}}N_j+C^{\mathrm{risk}}(C^{\mathrm{inc}}_j)}{v_j},
\end{equation}
with $C^{\mathrm{risk}}$ of Proposition~\ref{prop:tasks-security}; $c_j(U)$ has the same form with $F^U_j$, $K_j(U)$ and $\Omega_j(U)$.
\item[(b)] (Crossover.) Run and risk costs do not fall with $\pi$. They lower the crossover price by $\Delta c^{\mathrm{anc}}_j/[n_j-n_j(x^*)]$, where $\Delta c^{\mathrm{anc}}_j$ is the excess of these costs in $A$ over the binding alternative $x^*$ (if it stays binding), and automation never pays, at any inference price, unless $c^0_j(A)+c^{\mathrm{anc}}_j(A)<\underline c_j(0)$: a floor that cheaper inference cannot remove.
\item[(c)] (Volume.) Maintenance and revalidation multiply the volume threshold by the factor
\[
1+\delta_MT_A+\frac{C^{\mathrm{reval}}_jT_A}{T_{\mathrm{up}}F^A_j}.
\] The per-process incident hazard adds to the fixed cost; the per-unit hazard lowers the margin $\underline c_j-c^{A,\mathrm{var}}_j$.
\item[(d)] (Scope economies.) Technical debt lowers each edge's integration premium to $\eta^\times_e-c^{\mathrm{cpl}}N_iN_j/v_e$, and Propositions~\ref{prop:tasks-complement} and \ref{prop:tasks-aon} hold with this net premium. The process risk cost $C^{\mathrm{risk}}(\sum_{j\in\mathcal Z}C^{\mathrm{inc}}_j)$ is submodular in $\mathcal Z$, so it preserves supermodularity: technical debt is a diseconomy of integration scope and security an economy of it.
\end{enumerate}
\end{proposition}

\noindent\emph{Intuition and example.} Run costs behave like a second build, paid every year, and they put a floor under the cost of automation that falling inference prices never reach. With amortization over \icTA{} years, maintenance at \icDeltaM{} of the build a year and revalidation at \icRevalShare{} of the build every \icTup{} years multiply the volume threshold by \icVolFactor. In the invoice chain, full automation carries maintenance of \icMaint{} (\icDeltaM{} of \icFbuild{} of builds), revalidation of \icReval, technical debt of \icTechDebt{} (\icNint{} integrations at \icCint{} and three coupled pairs at \icCcpl) and, on \icCinc{} of expected incident loss, optimal security spending of \icSecSpend{} plus a residual loss of \icSecResid: \icRunTotal{} a year in all, engine defaults throughout. They turn the value of full automation at readiness \icIotaHigh{} in the conditional case from \icWfullHigh{} into \icWrunHigh, and no readiness level makes it positive.

\begin{proposition}[Build or buy]\label{prop:tasks-buildbuy}
Building costs $C^{\mathrm{fix}}_b:=F^A/T_A+\delta_MF^A+C^{\mathrm{reval}}/T_{\mathrm{up}}+C^{\mathrm{td}}$ a year plus $\pi n$ per unit. A vendor platform with build $F^V$, revalidation cost $C^{\mathrm{reval},V}$ and pooled volume $V_B$ charges $c_v:=(1+\upsilon_B)[\pi n+(F^V(1/T_A+\delta_M)+C^{\mathrm{reval},V}/T_{\mathrm{up}})/V_B]$ per unit, and buying costs an integration $F^{\mathrm{int}}$ and a lock-in premium $C^{\mathrm{lock}}$ a year, so $C^{\mathrm{fix}}_v:=F^{\mathrm{int}}/T_A+C^{\mathrm{lock}}$. Review and incident costs are the same either way, and $c_v>\pi n$.
\begin{enumerate}
\item[(a)] Building is cheaper if and only if $v\ge v^{BB}:=(C^{\mathrm{fix}}_b-C^{\mathrm{fix}}_v)/(c_v-\pi n)$.
\item[(b)] A higher maintenance rate favors buying if and only if $F^A>(1+\upsilon_B)(v/V_B)F^V$, that is, if the firm's own build exceeds its marked-up share of the vendor's; revalidation works the same way.
\item[(c)] (Lock-in.) Building wins if and only if $C^{\mathrm{lock}}\ge C^{\mathrm{lock}*}:=C^{\mathrm{fix}}_b-F^{\mathrm{int}}/T_A-v(c_v-\pi n)$. A vendor cannot raise its price by more than the switching cost $C^{\mathrm{sw}}$ per contract term $T_c$ without losing the client, so $C^{\mathrm{lock}}\le C^{\mathrm{sw}}/T_c$; if $C^{\mathrm{sw}}/T_c<C^{\mathrm{lock}*}$, buying is optimal whatever the lock-in.
\end{enumerate}
\end{proposition}

\noindent\emph{Intuition and example.} A vendor spreads maintenance and revalidation over all its clients, so run costs favor buying, and lock-in can offset this only up to the switching cost per contract term. For the invoice chain, building costs \icBBfix{} a year in fixed costs (\icBBbuilds{} of amortized builds and \icBBrun{} of run costs), while a platform with a \icBBmarkup{} markup, a pooled fixed cost of \icBBpooled{} per invoice and an integration of \icBBint{} costs \icBBfixv{} a year plus \icBBcv{} an invoice more than in-house inference. Building pays only above \icBBv{} invoices a year, against \icBBvnorun{} without run costs; at \icVolume{} invoices buying saves \icBBsave{} a year, above the \icBBlock{} a year that a \icBBsw{} switching cost on \icContractYears-year contracts would let a vendor extract.

\begin{proposition}[Security investment and scope]\label{prop:tasks-security}
Let $C^{\mathrm{inc}}_p:=\sum_{j\in\mathcal Z}C^{\mathrm{inc}}_j$ and $C^{\mathrm{risk}}(C^{\mathrm{inc}}_p):=\min_{C\ge 0}\{C+\mathsf g(C)C^{\mathrm{inc}}_p\}$.
\begin{enumerate}
\item[(a)] Optimal security spending $C^{\mathrm{sec}*}_p$ is positive if and only if $-\mathsf g'(0)C^{\mathrm{inc}}_p>1$, satisfies $-\mathsf g'(C^{\mathrm{sec}*}_p)C^{\mathrm{inc}}_p=1$ when positive, and is nondecreasing in $C^{\mathrm{inc}}_p$, hence in exposure and in the loss per incident. $C^{\mathrm{risk}}$ is nondecreasing, concave and subadditive, with slope $\mathsf g(C^{\mathrm{sec}*}_p)$: securing several activities or processes on one platform costs less than securing each separately.
\item[(b)] (Exponential residual.) If $\mathsf g(C)=\exp(-C/C_s)$, then $C^{\mathrm{sec}*}_p=C_s\ln(C^{\mathrm{inc}}_p/C_s)^+$, the residual expected loss at the optimum is $C_s$ whenever spending is positive, $C^{\mathrm{sec}*}_p\le C^{\mathrm{inc}}_p/e$ (the bound of \citealp{gordon2002economics}), and $C^{\mathrm{risk}}=C^{\mathrm{inc}}_p$ if $C^{\mathrm{inc}}_p\le C_s$ and $C_s[1+\ln(C^{\mathrm{inc}}_p/C_s)]$ otherwise.
\item[(c)] (Scope.) Activity $j$ is out of scope for automation if the firm's risk appetite excludes its data class, or if a single incident could exceed the loss $L^{\max}$ the firm tolerates. Otherwise it is worth automating only if its other savings exceed its marginal risk cost, which is at most $C^{\mathrm{risk}}(C^{\mathrm{inc}}_j)$; under (b), a stand-alone activity with other savings $W_j$ is in scope if and only if $W_j>C^{\mathrm{inc}}_j$ when $C^{\mathrm{inc}}_j\le C_s$ and $W_j>C_s[1+\ln(C^{\mathrm{inc}}_j/C_s)]$ otherwise. An activity excluded from $A$ can remain in scope for $U$, whose exposure is lower.
\end{enumerate}
\end{proposition}

\noindent\emph{Intuition and example.} Security spending is worth its marginal reduction in expected loss, so it grows with exposure, but less than proportionally, and the residual loss it leaves is set by how effective spending is, not by how large the exposure is. Once security is optimized, expected incident loss enters the scope test only logarithmically: \$500{,}000 of expected loss a year costs \icScope{} when spending \icCs{} cuts the hazard by a factor $e$ (the engine's default). What removes an activity from scope is therefore mostly risk appetite, a data class or a tail loss the firm will not bear, rather than expected loss. The bound in (b), that optimal spending never exceeds $1/e$ of the expected loss, is the result of \citet{gordon2002economics}, who prove it for two classes of breach-probability functions; the exponential residual here has the form of their second class, and the bound fails for other functions \citep{hausken2006returns}, so it should be read as a property of this specification rather than a general rule. Because the risk cost is concave, shared security platforms are an economy of scope that strengthens the case for integrating whole processes.

\begin{proposition}[Security review capacity]\label{prop:tasks-secreview}
Under Assumption~\ref{ass:tasks-ancillary}(iv), let each deployment be worth $V_{\mathrm{dep}}$ if it starts at once, with value eroding at the rate $r_d$ while it waits, as in Section~\ref{sec:ext-bandwidth}.
\begin{enumerate}
\item[(a)] Review latency $T_{\mathrm{sec}}$ is exponential with rate $\Phi_{\mathrm{sec}}-\nu$, so $\mathbb ET_{\mathrm{sec}}=1/(\Phi_{\mathrm{sec}}-\nu)$, and the expected value lost per deployment is $V_{\mathrm{dep}}r_d/(\Phi_{\mathrm{sec}}-\nu+r_d)$.
\item[(b)] The value lost per year is $D_{\mathrm{sec}}(\Phi_{\mathrm{sec}})=\nu V_{\mathrm{dep}}r_d/(\Phi_{\mathrm{sec}}-\nu+r_d)$, and a unit of extra review capacity is worth
\begin{equation}\label{eq:tasks-secvalue}
-D_{\mathrm{sec}}'(\Phi_{\mathrm{sec}})=\frac{\nu V_{\mathrm{dep}}r_d}{(\Phi_{\mathrm{sec}}-\nu+r_d)^2}\quad\text{a year.}
\end{equation}
\item[(c)] If capacity costs $p_{\mathrm{sec}}$ per review a year, the optimal capacity is $\Phi^*_{\mathrm{sec}}=\nu-r_d+\sqrt{\nu V_{\mathrm{dep}}r_d/p_{\mathrm{sec}}}$, provided this exceeds $\nu$: a square-root staffing rule.
\item[(d)] (Pace.) Deployments cannot be sustained faster than $\Phi_{\mathrm{sec}}$ a year. If deployments pass the decision gate of Section~\ref{sec:ext-bandwidth} and then security review, independently, the expected share of value retained is $\exp(-r_d\ell_{\mathrm{dec}})(\Phi_{\mathrm{sec}}-\nu)/(\Phi_{\mathrm{sec}}-\nu+r_d)$.
\end{enumerate}
\end{proposition}

\noindent\emph{Intuition and example.} Review capacity is a parallel bottleneck to decision bandwidth: it is cheap to add when the queue is short and very valuable when utilization is high, because latency explodes as $\nu$ approaches $\Phi_{\mathrm{sec}}$. With eight deployments a year worth \$1.5 million each, value eroding at 25\% a year and capacity for ten reviews, reviews take half a year on average and cost \$167{,}000 per deployment, \$1.33 million a year; a unit of extra capacity is worth \$593{,}000 a year, and at \$60{,}000 a unit the optimum is 14.8 reviews a year, which cuts latency to 0.15 years.

\begin{proposition}[How ancillary costs scale]\label{prop:tasks-ancscale}
\begin{enumerate}
\item[(a)] Ancillary costs split into a firm-level fixed part (compliance, governance, insurance), parts fixed per process (maintenance, revalidation, technical debt, security spending and the per-process hazard, security review, vendor and contract management, evaluation datasets) and parts per unit (per-transaction incident losses, inference operations). Per unit of work, the first two fall with the number of processes and with process volume.
\item[(b)] (Coarse data.) An analysis that treats $P$ processes as one cell counts per-process fixed costs once and so understates them by $\sum_p\phi_p-\bar\phi$; for equal processes, it understates total ancillary cost by the share $(1-1/P)s_{pp}$, where $s_{pp}$ is the share fixed per process. Imputing the mean integration count understates a convex technical-debt function by $\sum_pC^{\mathrm{td}}(N_p)-PC^{\mathrm{td}}(\bar N)\ge 0$, which is $(c/2)P\operatorname{Var}(N)$ for a quadratic with curvature $c$. If security is in fact bought process by process, pooling the cell's exposure understates the risk cost by $\sum_pC^{\mathrm{risk}}(C^{\mathrm{inc}}_p)-C^{\mathrm{risk}}(\sum_pC^{\mathrm{inc}}_p)\ge 0$.
\end{enumerate}
\end{proposition}

\noindent\emph{Intuition and example.} Ancillary costs are largely fixed per process, so coarse spend data, which cannot see processes, miss most of them. A cell of five processes like the invoice chain, each with \icScalePer{} of per-process maintenance and revalidation, integration counts from 2 to 10 and \icCinc{} of expected incident loss, is understated by \icScaleFixed{} of per-process fixed costs, \icScaleJensen{} of technical debt (at \icScaleCurv{} per coupled pair, an illustrative curvature) and \icScalePool{} of security cost if each process must be secured on its own. These are further reasons, beyond Proposition~\ref{prop:tasks-granularity}, to collect data at the process tier before committing to a roadmap.

\subsection{AI failure modes: hallucination, misbehavior and autonomy}\label{sec:tasks-failure}

The single raw error rate of Assumption~\ref{ass:tasks-review} hides three kinds of failure that controls treat differently.

\begin{assumption}[Failure taxonomy]\label{ass:tasks-failure}
In mode $A$ at autonomy level $\mathsf a\in\{0,1,2,3\}$ (read only, draft, act with approval, act autonomously), each unit of activity $j$:
\begin{enumerate}
\item[(i)] contains an ordinary error with probability $e^o_j$ and a hallucination (a confident fabrication) with probability $e^h_j$, which review of the unit catches with probabilities $d^o_j$ and $d^h_j<d^o_j$; $d^h_j$ is higher for retrieval-grounded than for open-ended tasks. Measured hallucination rates are 1.5\% to 20\% of outputs for retrieval-grounded work \citep{asgari2025framework,vectara2026hallucination}, 17\% to 34\% for legal research tools \citep{magesh2025hallucination}, and 5\% to 50\% for open generation of specific facts and citations \citep{walters2023fabrication,dahl2024large,openai2025o3}; catch rates, with the measured values widened for transfer, are 0.45 to 0.85 for ordinary errors and 0.2 to 0.6 for hallucinations \citep{panko1999applying,si2024large,dratsch2023automation}. Undetected errors cost $L_jz$, as in Assumption~\ref{ass:tasks-review};
\item[(ii)] triggers a misbehavior event (an out-of-scope action, tool misuse, specification gaming, following an injected instruction, concealment) with probability $e^m_j(\mathsf a)$, nondecreasing in autonomy and in the task horizon. The event is stopped before harm with probability $d^m_j(\mathsf a)$ (by approval at $\mathsf a\le 2$, by monitoring at $\mathsf a=3$); otherwise it costs $L^m_j(\mathsf a)\mathsf Z$, where the scale $L^m_j(\mathsf a)$ rises with write and transaction rights and $\mathsf Z\ge 1$ is Pareto with tail index $\mathsf p>0$, $P(\mathsf Z>x)=x^{-\mathsf p}$ (operational losses of banks, the class in which an AI execution failure would be booked, have tail indices of about 1.2 to 2.0, pooled 1.5, \citealp{defontnouvelle2007implications}, and IT cost overruns about 2.3, \citealp{flyvbjerg2022empirical}). Least privilege caps the loss at $L^{\mathrm{cap}}_j(\mathsf a)$, that is, $\mathsf Z$ at $\bar{\mathsf Z}_j(\mathsf a):=L^{\mathrm{cap}}_j(\mathsf a)/L^m_j(\mathsf a)$;
\item[(iii)] needs $t_j(\mathsf a)$ hours of human time, nonincreasing in $\mathsf a$.
\end{enumerate}
Injected-instruction events are counted here or in the incident cost of Section~\ref{sec:tasks-ancillary}, not in both. The firm's risk appetite bounds the conditional value at risk at level $\mathsf q$ of a single event's loss by $L^{\max}$.
\end{assumption}

\begin{proposition}[Errors and hallucination under review]\label{prop:tasks-halluc}
Let $\bar e_j:=e^o_j+e^h_j$ and $\tilde e_j:=e^o_jd^o_j+e^h_jd^h_j$, the catchable rate.
\begin{enumerate}
\item[(a)] With stakes-targeted review of the share $\xi$, the quality cost is $c^Q_j(\xi)=w_jt^R_j\xi+\tilde e_jL_j[1-\xi(1+\ln(1/\xi))]+(\bar e_j-\tilde e_j)L_j$. This is Proposition~\ref{prop:tasks-review} with raw error rate $\bar e_j$ and the error-weighted detection rate $d^R_j=\tilde e_j/\bar e_j$: the review-proof floor $(\bar e_j-\tilde e_j)L_j=(1-d^R_j)\bar e_jL_j$ is a loss that no review share removes.
\item[(b)] The undetected error rate is $\bar e_j-\tilde e_j\xi$. Review can meet the reliability requirement if and only if $e^o_j(1-d^o_j)+e^h_j(1-d^h_j)\le\bar\varepsilon_j$, the detection threshold of Proposition~\ref{prop:tasks-review}(d), and the reliability term of \eqref{eq:tasks-xi} becomes $(\bar e_j-\bar\varepsilon_j)/\tilde e_j$. Otherwise the activity is out of scope for automation until grounding raises $d^h_j$ or lowers $e^h_j$; raising $d^h_j$ by $\Delta$ lowers the floor by $e^h_j\Delta L_j$ per unit.
\end{enumerate}
\end{proposition}

\begin{proposition}[Automation bias: when more review buys less]\label{prop:tasks-bias}
Let units carry error probabilities $e_ju$, with a risk score $u\sim\mathrm{Exp}(1)$ observed before review, let review take the units with the highest scores, and let a reviewer catch an error with probability $d^R(\bar e)=d^R_0(\bar e/\bar e_0)^{\epsilon_d}<1$, where $\bar e$ is the error prevalence among reviewed units and $\epsilon_d>0$ the elasticity of vigilance: errors are missed more when they are rarer.
\begin{enumerate}
\item[(a)] The prevalence among reviewed units is $\bar e(\xi)=e_j(1+\ln(1/\xi))$, which falls as $\xi$ rises, and errors caught per unit are $\mathsf k(\xi)=e_j\xi(1+\ln(1/\xi))\,d^R(\bar e(\xi))$.
\item[(b)] (More review buys less.) $\mathsf k$ rises on $(0,\exp(-\epsilon_d))$ and falls on $(\exp(-\epsilon_d),1]$: reviewing more than the share $\exp(-\epsilon_d)$ catches fewer errors at higher cost, so the optimal review share never exceeds $\exp(-\epsilon_d)$.
\item[(c)] (Irony of automation.) $\partial(e_j-\mathsf k)/\partial e_j=1-(1+\epsilon_d)\mathsf k/e_j$: if review catches more than the share $1/(1+\epsilon_d)$ of errors, a more accurate model raises the undetected error rate.
\item[(d)] (Detection threshold.) The reliability target $\bar\varepsilon_j$ can be met at some review share if and only if
\[
d^R\bigl(e_j(1+\epsilon_d)\bigr)\ \ge\ \Bigl(1-\frac{\bar\varepsilon_j}{e_j}\Bigr)\frac{\exp(\epsilon_d)}{1+\epsilon_d}.
\]
Since $\exp(\epsilon_d)>1+\epsilon_d$, automation bias raises the detection rate that the target requires above the threshold $1-\bar\varepsilon_j/e_j$ of Proposition~\ref{prop:tasks-review}(d).
\end{enumerate}
\end{proposition}

\noindent\emph{Intuition and example.} Targeted review first takes the riskiest units, so each further unit reviewed dilutes the pool, and reviewers who see fewer errors find fewer of them. The direction of this effect is well documented: monitoring of automated systems deteriorates as they become more reliable, and training does not prevent automation bias \citep{parasuraman1993performance,bailey2007automation,parasuraman2010complacency}, while cognitive forcing reduces overreliance on AI \citep{bucinca2021trust}. Its size is not measured, so the following is conditional on the engine's elasticity of \icEpsd. Review beyond \icBiasPeak{} of units then catches fewer errors, and the detection rate that a reliability target requires is \icBiasUplift{} higher than without the bias. The same mechanism makes better models partly self-defeating: at a review share of one half, review catches \icBiasTop{} of errors when the catch rate is \icDRhigh, the top of the measured range, more than $1/(1+\epsilon_d)=\icBiasThr$, so a small gain in model accuracy raises the rate of errors that slip through; at the central measured rate of \icDR{} it catches \icBiasCentral, and the irony does not arise. In the invoice chain the review-proof floor binds, and ordinary errors rather than hallucinations set it. For approval at readiness \icIotaHigh, ordinary errors occur at \icHalEo{} and are caught at \icDR, and hallucinations at \icHalEh{} (the engine's retrieval-grounded rate of \icHret, lowered by grounding and capability) are caught at \icDh: the review-proof rate is \icHalFloor, of which \icHalFloorO{} comes from ordinary errors, against a tolerance of \icHalTol{} at every review share, and lifting the hallucination catch rate to 0.6 lowers it only to \icHalFloorDh. For open generation, with hallucination rates of 5\% or more, hallucinations alone put the floor at 3\% or more.

\begin{proposition}[Optimal autonomy]\label{prop:tasks-autonomy}
Under Assumption~\ref{ass:tasks-failure}, let
\[
c^m_j(\mathsf a):=e^m_j(\mathsf a)[1-d^m_j(\mathsf a)]L^m_j(\mathsf a)\,\mathbb E\min\{\mathsf Z,\bar{\mathsf Z}_j(\mathsf a)\}
\]
be the expected misbehavior loss per unit and $c_j(\mathsf a):=w_jt_j(\mathsf a)+c^Q_j(\mathsf a)+c^m_j(\mathsf a)$, other costs being equal across levels. The firm minimizes $c_j(\mathsf a)$ over the levels that satisfy its risk appetite.
\begin{enumerate}
\item[(a)] (Heavy tails.) $\mathbb E\min\{\mathsf Z,\bar{\mathsf Z}\}=1+(1-\bar{\mathsf Z}^{1-\mathsf p})/(\mathsf p-1)$, with $1+\ln\bar{\mathsf Z}$ at $\mathsf p=1$. Without a cap it is $\mathsf p/(\mathsf p-1)$ for $\mathsf p>1$ and infinite for $\mathsf p\le 1$: an infinite-mean tail rules out every level whose misbehavior is not always stopped, and least privilege restores a finite expected loss. Uncapped, the conditional value at risk of an event is $L^m_j\mathsf p(1-\mathsf q)^{-1/\mathsf p}/(\mathsf p-1)$ for $\mathsf p>1$; capped, it is at most $L^{\mathrm{cap}}_j$.
\item[(b)] (Optimal level.) Raising autonomy from $\mathsf a$ to $\mathsf a+1$ pays if and only if the human cost it saves, $w_j[t_j(\mathsf a)-t_j(\mathsf a+1)]$, exceeds the rise in quality and misbehavior cost, provided $\mathsf a+1$ satisfies the risk appetite.
\item[(c)] (Comparative statics.) If $c^m_j(\mathsf a)=\theta_m\hat c^m_j(\mathsf a)$ with $\hat c^m_j$ nondecreasing in $\mathsf a$, the largest optimal level is nonincreasing in $\theta_m$, which can stand for the misbehavior rate, the horizon factor, the loss scale or, with a common cap, the tail weight; it is nondecreasing in the wage.
\item[(d)] (Irreversible actions.) For an activity whose losses cannot be rolled back, mandatory human approval ($\mathsf a\le 2$) is required whenever the autonomous level violates the risk appetite, and the activity is out of scope for automation when even approved action does. It then stays augmented or in house.
\end{enumerate}
\end{proposition}

\noindent\emph{Intuition and example.} Autonomy trades human time for tail risk, and least privilege is what makes the trade possible when tails are heavy. Misbehavior of any kind occurs in 0.4\% to 1\% of autonomous tasks in evaluations of frontier agents \citep{vonarx2025recent,schoen2025stress}, and in about half of attempts when the scoring can be gamed \citep{zhong2025impossiblebench}; the consequential share, the effect of autonomy and the catch rate of human approval have not been measured, so the example is conditional on the engine's scenario values for them. For invoice approval the engine's inputs give a loss of \icAutL{} per undetected error, a misbehavior loss scale of \icAutLm{} after least privilege and monitoring, a kill-switch cap of \icAutCap{} and a tail index of \icAlpha. Expected misbehavior then costs \icAutCmTwo{} per invoice when actions are approved and \icAutCmThree{} when they are autonomous, against human time of \icAutHumOne{} per invoice to draft and \icAutHumTwo{} to approve. Uncapped, the conditional value at risk at \icAutCvarQ{} of an event is \icAutCvar{} at tail index \icAlpha{} and \icAutCvarLow{} at \icAlphaLow, far inside the engine's appetite of \icAutLmax. For an invoice-sized step misbehavior is not the binding risk; reliability is, and at the default detection rate neither automated level can meet approval's \icQbarAppr{} target, so the step stays augmented.

\begin{proposition}[Agent pipelines]\label{prop:tasks-pipeline}
Consider a pipeline of $J$ steps, each failing undetected with probability $\varepsilon$, independently. (Lemma~\ref{lem:tasks-oring} is the case without checkpoints and with losses that do not cascade.)
\begin{enumerate}
\item[(a)] The pipeline succeeds with probability $(1-\varepsilon)^J$, and on $k$ independent tasks in a row (pass$^k$) with probability $(1-\varepsilon)^{Jk}$. The maximum length that meets a target $\bar q_p$ is $J_{\max}=\lfloor\ln\bar q_p/[k\ln(1-\varepsilon)]\rfloor$, about $(1-\bar q_p)/(k\varepsilon)$.
\item[(b)] (Checkpoints.) A checkpoint after every $J_s$ steps that catches a failed segment with probability $d^c$ leaves an undetected failure probability $[1-(1-\varepsilon)^{J_s}](1-d^c)$ per segment, and the maximum length becomes $J_s\ln\bar q_p/\ln\{1-[1-(1-\varepsilon)^{J_s}](1-d^c)\}$, about $J_{\max}/(1-d^c)$ for small $\varepsilon$.
\item[(c)] (Cascading losses.) If an undetected error at step $i$ is acted on by every later step, each adding a loss $L^c$, the expected loss per run is $L^c[J-(1-\varepsilon)(1-(1-\varepsilon)^J)/\varepsilon]=L^c\varepsilon J(J+1)/2+O(\varepsilon^2)$: it grows with the square of the pipeline's length, and the longest pipeline within a loss budget $\bar L$ per run is about $\sqrt{2\bar L/(\varepsilon L^c)}$.
\item[(d)] (Decomposition.) If step $k$ takes $d_k$ hours of human time and succeeds with probability $2^{-d_k/\mathcal H}$ (a constant hazard), the pipeline succeeds with probability $2^{-\sum_kd_k/\mathcal H}$, which depends on its total human time and not on its number of steps. Under the logistic form of Assumption~\ref{ass:ext-capability} with $\eta_r\ge 1$, splitting any step into two independent steps of the same total length lowers the pipeline's success probability.
\end{enumerate}
\end{proposition}

\noindent\emph{Intuition and example.} Reliability requirements compound along a pipeline and over repeated tasks: at 1\% per step, a pipeline that must succeed 95\% of the time can have five steps, and only one if it must succeed on five tasks in a row. Consistency across repeated trials is the binding statistic for high-volume work: an agent that succeeds on 61\% of single trials succeeds on all of eight trials less than 25\% of the time \citep{yao2025tau}, and multi-agent systems fail on 41\% to 87\% of tasks, a third of the failures arising between agents \citep{cemri2025why}. On steps of an hour of human time, per-step reliability is 0.74 to 0.94, the two ends being the logistic and constant-hazard readings of the horizon data \citep{kwa2025measuring,ord2025is}: no such step meets a 95\% target alone, and three steps at 0.94 meet only 80\%. Checkpoints stretch the feasible length by $1/(1-d^c)$, a fivefold gain if a checkpoint catches 80\% of failed segments (a scenario value), while cascading losses shorten it, because an early error contaminates everything downstream. Part (d) says that decomposition alone buys no reliability; it pays only through the checkpoints between steps.

\begin{proposition}[The error and misbehavior loss floor]\label{prop:tasks-cerr}
(a) After optimal review and autonomy, the expected loss per unit from errors, hallucinations and misbehavior is
\[
c^{\mathrm{err}}_j:=\tilde e_jL_j[1-\xi^*_j(1+\ln(1/\xi^*_j))]+(\bar e_j-\tilde e_j)L_j+c^m_j(\mathsf a^*_j)\ \ge\ (\bar e_j-\tilde e_j)L_j+c^m_j(\mathsf a^*_j),
\]
and the lower bound, which no review share removes, falls only with grounding, least privilege and lower failure rates. (b) As a share of revenue at full automation,
\begin{equation}\label{eq:tasks-cerr}
c_{\mathrm{err}}:=\frac1S\sum_{j\in\mathcal J^A}v_jc^{\mathrm{err}}_j ,
\end{equation}
where $\mathcal J^A$ is the set of automated activities. It enters Part I's full-automation ceiling as a cost alongside $c_c$, $f$ and $a_{\min}$, so that $(1-\tau)(1-c_c-f-a_{\min}-c_{\mathrm{err}}-\cdots)$ replaces $(1-\tau)(1-c_c-f-a_{\min})$ in Corollary~\ref{cor:ceiling}. Review labor is excluded from $c_{\mathrm{err}}$ (it is human cost), and so are the security incidents of Section~\ref{sec:tasks-ancillary}. (c) (Full automation without people.) If no person reviews, the catch rates in (a) are those of the review that remains. With no review, $\xi^*_j=0$ and $c^{\mathrm{err}}_j=\bar e_jL_j+c^m_j(\mathsf a^*_j)$. With review by a model that shares the producing model's error with probability $\rho^{\mathrm{cm}}_j$ (common mode) and otherwise catches it with probability $d^A_j$, the catch rate of an error is $(1-\rho^{\mathrm{cm}}_j)d^A_j$, which replaces the human rates in (a); $\rho^{\mathrm{cm}}_j$ is close to one for a reviewer on the same model, and for one on another provider it is the probability that the second model also errs and agrees, of which only the agreement given a joint error has been measured. Evaluating $c_{\mathrm{err}}$ at human catch rates for a firm without people therefore understates it.
\end{proposition}

\begin{corollary}[Failure and the success of a transition]\label{cor:tasks-failure-success}
Suppose a transition succeeds only if the deployed process meets its reliability target, $\prod_{j\in\mathcal C_p}(1-\varepsilon_j)\ge\bar q_p$, with uncertain step rates $\varepsilon_j$. Then $\varpi_p$ is zero if some critical-path step has a review-proof rate $e^o_j(1-d^o_j)+e^h_j(1-d^h_j)$ above $1-\bar q_p$; it falls with the number of steps and with each step's hallucination and misbehavior rates and its autonomy level; and it rises with checkpoints. The log reliability margin $\ln(J_{\max}/J)$ is therefore a natural driver in the logit of the factors $\varpi_{pj}$ in \eqref{eq:tasks-success}.
\end{corollary}

\noindent\emph{The floor in numbers.} For invoice approval at readiness \icIotaHigh, with the error and catch rates of the hallucination example and a loss of \icAutL{} per undetected error, the floor is \icCerrFloor{} of review-proof errors plus \icCerrMis{} of misbehavior (conditional on the scenario values of the autonomy example) per invoice, \icCerrAnnual{} a year on \icVolume{} invoices: a cost that persists however cheap inference becomes. That floor assumes people review; with no review at all, part (c), it rises to \icCerrNoRev{} per invoice, \icCerrNoRevAnnual{} a year, and review by a model lies between the two by how much of the producing model's error it shares.

\subsection{The value of granular spend data}\label{sec:tasks-granularity}

Data arrive in tiers, from public financials (Tier 0) to process event logs (Tier 4). Model tier $T$ as information $\mathcal I_T$, with $\mathcal I_0\subseteq\mathcal I_1\subseteq\cdots$. Let $Y_j$ be the annual value of automating activity $j$ against its best alternative, uncertain because finer levels are imputed. The firm automates $j$ if and only if $\mathbb E[Y_j\mid\mathcal I_T]>0$, and its expected value at tier $T$ is $\mathcal V_T:=\sum_j\mathbb E[(\mathbb E[Y_j\mid\mathcal I_T])^+]$.

\begin{proposition}[The value of granularity]\label{prop:tasks-granularity}
Assume $\mathbb E|Y_j|<\infty$.
\begin{enumerate}
\item[(a)] $\mathcal V_T$ is nondecreasing in $T$. The expected annual loss from misassigned modes, $\mathcal L_T:=\sum_j\mathbb E[Y_j^+]-\mathcal V_T$, is nonnegative, and moving from tier $T$ to $T'>T$ is worth $\mathcal V_{T'}-\mathcal V_T\ge 0$.
\item[(b)] (Distance to the crossover.) Let $m_j$ and $s_j$ be the mean and standard deviation of $Y_j$ given $\mathcal I_T$. The loss for activity $j$, $\ell_j:=\mathbb E[Y_j^+\mid\mathcal I_T]-m_j^+$, satisfies $\ell_j\le\min\{s_j/2,\,s_j^2/(4|m_j|)\}$ for every posterior distribution.
\item[(c)] (Normal posteriors.) If $Y_j\mid\mathcal I_T\sim N(m_j,s_j^2)$, then $\ell_j=s_j\mathsf L(|m_j|/s_j)$ with $\mathsf L(u):=\mathsf n(u)-u[1-\mathsf N(u)]\le\mathsf n(u)/(1+u^2)$, where $\mathsf N$ and $\mathsf n$ are the standard normal distribution and density; $\mathsf L$ decreases from $\mathsf L(0)\approx 0.399$. A finer tier that lowers the posterior variance to $s_j'^2$ is worth $\hat s_j\mathsf L(|m_j|/\hat s_j)$, with $\hat s_j^2=s_j^2-s_j'^2$.
\item[(d)] (What spend data do.) Write $Y_j=E_j\delta_j-F^A_j/T_A$, with spend $E_j$ and saving rate $\delta_j$ independent given $\mathcal I_T$. If $F^A_j=0$, observing $E_j$ has no value for the mode decision. If $F^A_j>0$ and $m_\delta:=\mathbb E[\delta_j\mid\mathcal I_T]>0$, observing it is worth $m_\delta\{\mathbb E[(E_j-E^*_j)^+]-(\mathbb E E_j-E^*_j)^+\}$, with spend threshold $E^*_j=F^A_j/(T_Am_\delta)$.
\item[(e)] (Cell-level action.) If the firm must choose one mode for a cell of activities it cannot tell apart, its loss against activity-level information is $\min\{\sum_jY_j^+,\sum_jY_j^-\}$ over the cell.
\end{enumerate}
\end{proposition}

\noindent\emph{Intuition.} Information has value only if it can change a decision \citep{howard1966information}, so coarse data suffice for activities whose estimated value lies far from a crossover, measured in posterior standard deviations. An activity with an estimated saving of \$200{,}000 a year and a standard deviation of \$300{,}000 loses \$45{,}000 a year in expectation from deciding on coarse data (the distribution-free bound is \$112{,}500); with an estimate of \$1 million the expected loss is \$34 (bound \$22{,}500). Part (d) says that finer spend data matter for the mode decision mainly near volume thresholds; what they improve elsewhere is the estimate of the prize, not the choice. Granularity stops paying where cells become homogeneous in their optimal mode, which part (e) measures, or where the value of the next tier falls below its cost.

\subsection{Aggregation into the frontier}\label{sec:tasks-aggregation}

Group activities into cost of revenue ($\mathcal J^{\mathrm{cogs}}$), non-discretionary operating cost ($\mathcal J^{\mathrm{opex}}$: general and administrative cost and maintenance R\&D) and growth activities ($\mathcal J^{\mathrm{grow}}$: selling, marketing and expansion R\&D). Growth spend buys a fixed bundle of growth activities. Split each unit cost into labor and services $c^{\mathrm{lab}}_j$ (in-house and review labor, vendor prices, coordination and error losses), inference $\pi n_j$, the amortized build and the capital charge; the first two make up the operating unit cost $c^{\mathrm{op}}_j$.

\begin{proposition}[Aggregation]\label{prop:tasks-aggregate}
Let the firm move from modes $x^0$ to $x^1$, and write $\Delta$ for the change, new minus old.
\begin{enumerate}
\item[(a)] (Mapping.) $\Delta c_h=\sum_{\mathcal J^{\mathrm{cogs}}}v_j\Delta c^{\mathrm{lab}}_j/S$, $\Delta c_c=\sum_{\mathcal J^{\mathrm{cogs}}}v_j\pi\Delta n_j/S$, $\Delta a=\sum_{\mathcal J^{\mathrm{opex}}}v_j\Delta c^{\mathrm{op}}_j/S+\sum_j\Delta(F^A_j/T_A)/S$ (builds are fixed costs, Remark~\ref{rem:assumptions}(a)), $\Delta k=\sum_jv_j\Delta K_j/S$, and $\lambda^1=\lambda^0(1-\delta_\sigma)^{-\alpha}$, where $\delta_\sigma:=-\sum_{\mathcal J^{\mathrm{grow}}}v_j\Delta c^{\mathrm{op}}_j/\sum_{\mathcal J^{\mathrm{grow}}}v_jc^{\mathrm{op}}_j(x^0_j)$ is the proportional saving on growth activities. With pass-through $\varphi$ (Lemma~\ref{lem:prices}), $\Delta\gamma=-(1-\varphi)(\Delta c_h+\Delta c_c)$ for a saving the firm alone makes. If every rival makes the same savings, the industry passes the rate $\bar\varphi^{\rm ind}$ of every unit-cost saving to the price, non-discretionary cost included (Part~I's mapping, which exposes all savings to competition), and the margin is computed at the lower price: with $C_0:=1-f-m_0$ today's non-fee unit cost per dollar of price and $S_u:=-(\Delta c_h+\Delta c_c+\Delta a)$ the saving, $m_0$ becomes $1-f-(C_0-S_u)/(1-\bar\varphi^{\rm ind}S_u)$, which is $m_0+(1-\bar\varphi^{\rm ind}C_0)S_u$ to first order; the elasticity of Lemma~\ref{lem:prices} is then $\varphi=\bar\varphi^{\rm ind}C_0$, and the concentration of the market does not enter. The part of the accountability floor that the activities represent is $a_{\min}=\sum_{\mathcal J^{\mathrm{opex}}}v_jc^{\mathrm{floor}}_j/S$, with $c^{\mathrm{floor}}_j:=\min_{x\in\mathcal M_j}[c^{\mathrm{op}}_j(x)+\mathbf 1\{x=A\}F^A_j/(v_jT_A)]$ evaluated at $\pi=0$ and $e_j=0$. This is the cost part of the floor. Proposition~\ref{prop:acct-amin} adds a condition, not a loss: the control design over each in-scope objective must leave no material weakness, which can rule out cheap all-AI designs even at $e_j=0$ when their layers share common-mode failures, and $a_{\min}$ is the cost of the cheapest design that passes. The residual expected consequence of misstatements the chosen design misses is an operating cost, not part of $a_{\min}$; the error losses $L_j$ and the floor $c_{\mathrm{err}}$ of Proposition~\ref{prop:tasks-cerr} exclude it, so it is counted once. Run costs, technical debt, security spending and firm-level compliance costs enter $a$ as fixed costs, and expected incident losses enter the bucket of the activity that causes them. Changes in retention growth $G_i$ from customer-facing activities are an input.
\item[(b)] (Score.) The new maximal score is $R_X^*$ of \eqref{eq:Rstar} at the new parameters, with $m_0=\gamma-a-\delta_Kk$. To first order, $\Delta R_X^*=-(1-\varphi)(\Delta c_h+\Delta c_c)-\Delta a-(g^*+\delta_K)\Delta k+\sigma^*\delta_\sigma+(X-k)\Delta G_i$, and $-(1-\varphi)(\Delta c_h+\Delta c_c+\Delta a)$ replaces the first two terms, with $\varphi=\bar\varphi^{\rm ind}C_0$, when every rival makes the savings. The exact effect of $\delta_\sigma$ alone is $\Pi[(1-\delta_\sigma)^{-\alpha/(1-\alpha)}-1]\ge\sigma^*\delta_\sigma$.
\item[(c)] (Firm value.) At small share the revenue multiple $v_0$ of \eqref{eq:gordon} satisfies $rv_0=R^*_X$ at $X=v_0$: the multiple is the maximal score at the firm's own weight, capitalized at $r$. With $X^\circ:=v_0$, to first order $\Delta v_0=\Delta R^*_{X^\circ}/(r-g^\circ)$, where $\Delta R^*_{X^\circ}$ is the change at fixed weight, and firm value changes by $S\Delta v_0$ less the one-time transition costs of Proposition~\ref{prop:tasks-transition}.
\end{enumerate}
\end{proposition}

\begin{corollary}[What a saved dollar is worth]\label{cor:tasks-weights}
To first order, a saving of one dollar per dollar of revenue raises $R_X^*$ by 1 in non-discretionary operating cost, by $1-\varphi$ in cost of revenue and by 1 in growth activities (more than 1 exactly, because spend is re-optimized), and by $1-\bar\varphi^{\rm ind}C_0$ in either kind of operating cost if every rival makes it; a dollar of capital per dollar of revenue lowers it by $g^*+\delta_K$; and a point of retention growth raises it by $X-k$. Each effect on the revenue multiple is the effect on $R^*_{X^\circ}$ divided by $r-g^\circ$.
\end{corollary}

\noindent\emph{Intuition and example.} The mapping makes each activity-level decision visible in the frontier: cost-of-revenue savings count only to the extent they are not passed through, savings on growth activities buy the same growth for less, and capital enters at the growth rate plus depreciation. Part (c) turns the score into value: since $rv_0=R^*_{v_0}$, a permanent gain in the score at the firm's own weight is capitalized at $1/(r-g^\circ)$. In Part I's worked example ($v_0=10$, $g^\circ=7.6\%$, $r=0.10$), saving one point of revenue in general and administrative cost raises the revenue multiple by 0.41 to first order (0.42 exactly), and cutting the cost of growth activities by 20\% adds 2.1 points to $R^*_{X^\circ}$ to first order (2.5 exactly) and raises the multiple from 10 to 11.3. These values are private. Where wages exceed workers' outside options, part of the saving is a transfer of rents, which \citet{acemoglu2026automation} estimate offsets 60 to 90\% of automation's measured productivity gains in the aggregate.

\subsection{Breakpoints}\label{sec:tasks-breakpoints}

\begin{proposition}[Breakpoints]\label{prop:tasks-breakpoints}
Fix all parameters but one. Automation of activity $j$ is optimal if and only if $c_j(A)\le\underline c_j$, and solving this one inequality for one variable at a time gives:
\begin{enumerate}
\item[(1)] (Cost crossover.) $\pi\le\pi^*_j$ (Proposition~\ref{prop:tasks-modes}(d)).
\item[(2)] (Volume.) $v_j\ge v^*_j$ (Proposition~\ref{prop:tasks-modes}(e)).
\item[(3)] (Reliability.) $e_j\le e^*_j$, where $e^*_j$ solves $c^Q_j(e)=\underline c_j-\pi n_j-F^A_j/(v_jT_A)-(r+\delta_K)K_j(A)$ when the right side lies between $\xi^{\min}_jw_jt^R_j$ and $c^Q_j$ at $e=1$; automation is infeasible at any review share if $d^R_j<1-\bar\varepsilon_j/e_j$, and as $e_j\to 0$ the review share falls to $\xi^{\min}_j$ (Proposition~\ref{prop:tasks-review}(d)).
\item[(4)] (Readiness.) $\iota_j\ge\iota^*_j$ (Proposition~\ref{prop:tasks-augment}), and for a process $\iota\ge\iota^*_p$ (Proposition~\ref{prop:tasks-readiness}).
\item[(5)] (Percolation.) Let activity $j$ be feasible for automation if and only if effective capability $\mathcal H$ is at least $\bar d_j$ (the test of Assumption~\ref{ass:ext-capability}), and let $\mathcal C_p=\mathcal J_p$. Under Proposition~\ref{prop:tasks-complement}(a), the largest optimal automated set of process $p$ is nondecreasing in $\mathcal H$ and changes only when $\mathcal H$ crosses some $\bar d_j$. If every proper subset of $\mathcal J_p$ has nonpositive value and $W_p(\mathcal J_p)>0$, the process flips from no automation to full automation exactly when $\mathcal H$ reaches $\bar d_p:=\max_{j\in\mathcal J_p}\bar d_j$.
\item[(6)] (Transition.) The process moves if and only if its horizon exceeds $T_h^*$ (Proposition~\ref{prop:tasks-transition}(b)).
\item[(7)] (Run-cost floor and scope.) Automation can pay at some inference price only if $c^0_j(A)+c^{\mathrm{anc}}_j(A)<\underline c_j(0)$ (Proposition~\ref{prop:tasks-fullcost}(b)), and only for activities in scope (Proposition~\ref{prop:tasks-security}(c)).
\item[(8)] (Review capacity.) Deployments cannot be sustained faster than $\Phi_{\mathrm{sec}}$ a year (Proposition~\ref{prop:tasks-secreview}(d)).
\item[(9)] (Failure.) Automation requires a review-proof error rate within tolerance (Proposition~\ref{prop:tasks-halluc}(b)), an autonomy level within the risk appetite (Proposition~\ref{prop:tasks-autonomy}(d)) and a pipeline no longer than $J_{\max}$ (Proposition~\ref{prop:tasks-pipeline}).
\end{enumerate}
Consequently, the automated share is a step function of the inference price, readiness and capability; under the condition of (5), it jumps by a whole process's spend at each $\bar d_p$.
\end{proposition}

\noindent\emph{Intuition.} The five activity and process thresholds are one inequality seen from different sides, so they move together: the crossover price rises with volume, readiness and the detection rate of review and falls with the loss per error, the reliability requirement and the oversight floor. Because processes automate whole once their hardest activity becomes feasible, steady growth in capability produces automation in process-sized lumps. The engine reports, for each activity and process, its distance to every breakpoint in natural units and in posterior standard deviations, which by Proposition~\ref{prop:tasks-granularity}(b) is also a measure of how much better data would be worth.
